% BEGIN INCLUDED SGA7I_English_current_buildable_source_20260823/components/09_expose_VI_opening_through_1_1.tex
%%%%%%%%%%  scan idx 43  |  folio 32  %%%%%%%%%%

\begin{center}
\underline{\SGARefAliasTarget{sga7.sga7i.evi.nroot}报告 VI}

\vspace{1em}
\CJKunderdblline{形式形变理论}

\vspace{1em}
\underline{D. S. Rim 著}
\end{center}

1. \underline{形式存在定理}

回顾余纤维范畴(cofibered category)的概念(S.G.A. 1 IV). 设
$\underline C$ 为固定范畴. 按定义, $\underline C$ 上的余纤维范畴
$\underline A$ 是配有函子 $P:\underline A\longrightarrow\underline C$
的范畴 $\underline A$, 满足:

\underline{公理 1.} 对 $\underline C$ 中每个映射 $f$ 及 $\underline A$
中每个满足 $P(a)$ 为 $f$ 的源的对象 $a$, 存在余笛卡尔的 $f$-态射
$\alpha:a\longrightarrow b$. 即对每个 $f$-态射 $\alpha':a\longrightarrow b'$,
存在 $\underline A$ 中唯一的 $T$-态射 $\tau:b\longrightarrow b'$
使 $\alpha'=\tau\alpha$, 其中 $T$ 为 $f$ 的靶.

\underline{公理 2.} $\underline A$ 中余笛卡尔态射的复合仍是余笛卡尔态射.

余纤维范畴 $P:\underline A\longrightarrow\underline C$ 的下述性质直接由定义得到.

(消去性) 设 $a\overset{\alpha}{\underset{\beta}{\rightrightarrows}}a'$
为两个余笛卡尔映射, 且 $P(\alpha)=P(\beta)$. 若存在余笛卡尔映射
$b\xrightarrow{\gamma}a$ 使 $\alpha\gamma=\beta\gamma$, 则 $\alpha=\beta$.

(分解性) 给定余笛卡尔映射 $a\xrightarrow{\alpha}a'$ 和
$a\xrightarrow{\beta}a''$, 存在余笛卡尔映射 $a'\xrightarrow{\gamma}a''$
使 $\gamma\alpha=\beta$, 当且仅当 $\underline C$ 中存在映射
$f:P(a')\longrightarrow P(a'')$ 使 $fP(\alpha)=P(\beta)$.
而且 $\gamma$ 唯一.

设 $\underline A$ 为 $\underline C$ 上的余纤维范畴. 对 $\underline C$
的每个对象 $S$, 用 $\underline A(S)$ 表示 $\underline A$ 的子范畴,
其对象是满足 $P(a)=S$ 的对象 $a$, 态射为
\[
 \Hom_{\underline A(S)}(a,a')=
 \{\alpha\in\Hom_{\underline A}(a,a')\mid P(\alpha)=\id_S\}.
\]
$\underline C$ 上的余纤维群胚(cofibered groupoid)是 $\underline C$
上的余纤维范畴 $\underline A$, 满足对 $\underline C$ 的每个对象 $S$,

%%%%%%%%%%  scan idx 44  |  folio 33  %%%%%%%%%%

$\underline A(S)$ 都是群胚. 回顾群胚是每个态射都是同构的范畴.
因此, $\underline C$ 上的余纤维群胚就是 $\underline A$ 中每个映射都余笛卡尔的余纤维范畴
$P:\underline A\longrightarrow\underline C$. 此外, 对每个
$\alpha\in\operatorname{Fl}(\underline A)$, $\alpha$ 是 $\underline A$
中的同构当且仅当 $P(\alpha)$ 是同构.

本节余下部分固定一个剩余域为 $k$ 的完备 诺特 局部环 $\Lambda$,
并以 $\underline C_\Lambda$ 表示剩余域为 $k$ 的 Artin 局部
$\Lambda$-代数所成范畴, 其映射为 $\Lambda$ 上的局部同态.
从 \SGARefLink{sga7.sga7i.evi.n1-2}{1.2} 起, 总假定
$\underline C_\Lambda$ 上的余纤维群胚 $\underline A$ 满足
$\underline A(k)=\{e\}$, 其中 $\{e\}$ 表示这样的范畴: 对 $\{e\}$ 的任意两个对象
$a,b$, $\Hom_{\{e\}}(a,b)$ 恰有一个元素.
$\underline A$ 中的映射 $a'\xrightarrow{\alpha}a$, 若满足
$P(a')=P(a)$ 且 $P(\alpha)=1_{P(a)}$, 就简称为
$\underline C_\Lambda$-映射. 稍微滥用术语, 若
$P(\alpha):P(a')\longrightarrow P(a)$ 在 $\underline C_\Lambda$
中满射, 就称 $\underline A$ 中的映射 $a'\xrightarrow{\alpha}a$ 为满射.

\underline{例 \SGARefTarget{sga7.sga7i.evi.n1-1.item-a}{\SGARefAliasTarget{sga7.sga7i.evi.n1-1}1.1}(a).}
设 $F:\underline C\longrightarrow(\text{范畴})$ 为函子.
定义范畴 $\underline F$ 如下. $\underline F$ 的对象为二元组 $(S,a)$,
其中 $S\in\mathrm{ob}(\underline C)$ 且 $a\in\mathrm{ob}F(S)$;
映射 $(S,a)\longrightarrow(S',a')$ 为二元组 $(f,u')$, 其中
$f:S\longrightarrow S'$ 是 $\underline C$ 中的映射,
$u':F(f)a\longrightarrow a'$ 是 $F(S')$ 中的映射.
\[
 (S,a)\xrightarrow{(f,u')}(S',a')
 \xrightarrow{(g,u'')}(S'',a'')
\]
的复合为
\[
 (gf,\,u''\mathbin{\cdot}F(g)u'):(S,a)\longrightarrow(S'',a'').
\]
对于由 $(S,a)\longmapsto S$ 给出的函子
$P:\underline F\longrightarrow\underline C$, $\underline F$
成为 $\underline C$ 上的余纤维范畴. 若对每个
$S\in\mathrm{ob}(\underline C)$, $F(S)$ 均为群胚,
则 $\underline F$ 是 $\underline C$ 上的余纤维群胚.
特别地, 函子 $F:\underline C\longrightarrow(\text{群})$
给出一个余纤维群, 函子 $F:\underline C\longrightarrow(\mathrm{Ens})$
给出一个余纤维离散群胚. 自然变换 $\varphi:F\longrightarrow G$,
其中 $F,G:\underline C\longrightarrow(\text{范畴})$, 诱导
$\underline C$ 上的余笛卡尔函子
$\varphi:\underline F\longrightarrow\underline G$.

\underline{例 \SGARefTarget{sga7.sga7i.evi.n1-1.item-b}{1.1}(b).}
设 $X$ 和 $Y$ 为 $\Lambda$ 上的概形, 并给定同构

%%%%%%%%%%  scan idx 45  |  folio 34  %%%%%%%%%%

\[
 X\otimes_\Lambda k\xleftrightarrow{\ \tau\ }Y\otimes_\Lambda k.
\]
则由
\[
 \Isom_{X,Y}(S)=
 \left\{X\otimes_\Lambda S\xleftrightarrow{\ \sigma\ }
 Y\otimes_\Lambda S\;\middle|\;
 \sigma\otimes_\Lambda k=\tau\right\}
\]
给出的函子 $\Isom_{X,Y}:\underline C_\Lambda\longrightarrow(\text{群胚})$
给出余纤维群胚 $P:\underline{\Isom}_{X,Y}\longrightarrow\underline C_\Lambda$.

\underline{例 \SGARefTarget{sga7.sga7i.evi.n1-1.item-c}{1.1}(c).}
设 $X_0$ 为 $k$ 上的固定概形, $\underline M_{X_0}$ 为
$X_0$ 在 $\underline C_\Lambda$ 上的形变范畴(见第 4 节).
因此 $\underline M_{X_0}$ 的对象为二元组 $(R,X)$,
其中 $R\in\underline C_\Lambda$, 而 $X$ 为 $X_0$ 到 $R$ 的形变.
于是由 $(R,X)\longmapsto R$ 给出的函子
$P:\underline M_{X_0}\longrightarrow\underline C_\Lambda$
定义 $\underline C_\Lambda$ 上的余纤维群胚.

\underline{例 \SGARefTarget{sga7.sga7i.evi.n1-1.item-d}{1.1}(d).}
对任意群胚 $X$, 以 $[X]$ 表示由 $X$ 的对象的同构类组成的离散群胚.
若 $P:\underline F\longrightarrow\underline C_\Lambda$ 是余纤维群胚,
则得到函子 $[\underline F]:\underline C_\Lambda\longrightarrow(\mathrm{Ens})$,
由 $[\underline F](S)=[\underline F(S)]$ 给出,
因而得到一个余纤维离散群胚; 稍微滥用记号, 仍将其记为 $[\underline F]$.

% END INCLUDED SGA7I_English_current_buildable_source_20260823/components/09_expose_VI_opening_through_1_1.tex

% BEGIN INCLUDED SGA7I_English_current_buildable_source_20260823/components/10_expose_VI_1_2_through_1_3.tex
\underline{定义 \SGARefTarget{sga7.sga7i.evi.n1-2}{1.2}.
称 $\underline C_\Lambda$ 上的余纤维群胚 $\underline A$ 为半齐性的(semi-homogeneous),
若满足下列性质:}

(1) \underline{$\underline A$ 中每个由实线箭头组成的图}
\[
\begin{tikzcd}[row sep=large,column sep=large]
a'' \arrow[r,dashed] \arrow[d,dotted] & a' \arrow[d] \\
a_1 \arrow[r] & a
\end{tikzcd}
\]
\underline{只要 $a'\longrightarrow a$ 满射, 就能添加点线箭头将其补为交换图.}

(2) \underline{设}
\[
\begin{tikzcd}[column sep=large]
& R \\
R\underset{k}{\times}k[\varepsilon]
 \arrow[ur,"\pi_1"] \arrow[dr,"\pi_2"'] & \\
& k[\varepsilon]
\end{tikzcd}
\]
\underline{为投影映射, 其中 $k[\varepsilon]$ 是 $k$ 上的对偶数环, 并设}
\[
\begin{tikzcd}[column sep=huge]
a' \arrow[dr,"\alpha"] & \\
& a \\
b' \arrow[ur,"\beta"'] &
\end{tikzcd}
\]
\underline{为 $\underline A$ 中的 $\pi_1$-映射.
则存在 $\underline C_\Lambda$-映射 $\gamma:a'\longrightarrow b'$
使 $\alpha=\beta\gamma$, 当且仅当
$(\pi_2)_*(a')\simeq(\pi_2)_*(b')$.}

\underline{注 \SGARefTarget{sga7.sga7i.evi.n1-3.item-a}{\SGARefAliasTarget{sga7.sga7i.evi.n1-3}1.3}(a).}
稍后将在 \SGARefLink{sga7.sga7i.evi.n2-6.item-b}{2.6}(b) 用 $2$-范畴重新表述这个概念.
若 $\underline A$ 半齐性, 则 $[\underline A]$ 也半齐性.
特别地, \SGARefLink{sga7.sga7i.evi.n1-2}{1.2}(2) 蕴含典范映射

%%%%%%%%%%  scan idx 46  |  folio 35  %%%%%%%%%%

\[
 [\underline A(R\times_k k[\varepsilon])]
 \longrightarrow
 [\underline A(R)]\times[\underline A(k[\varepsilon])]
\]
为双射.

\underline{注 \SGARefTarget{sga7.sga7i.evi.n1-3.item-b}{1.3}(b).}
考虑 $\underline C_\Lambda$ 上余纤维范畴 $\underline A$ 中的交换图
\[
\begin{tikzcd}[row sep=large,column sep=large]
& a' \arrow[dl,"v_1"'] \arrow[dr,"v_2"] & \\
a_1 \arrow[dr,"u_1"'] & & a_2 \arrow[dl,"u_2"] \\
& a &
\end{tikzcd}
\]
置 $R=P(a)$, $R'=P(a')$, 且对 $i=1,2$ 置 $R_i=P(a_i)$.
若 $a''$ 为映射
\[
 (P(v_1),P(v_2)):R'\longrightarrow R_1\times_RR_2,
\]
下的正像, 则得到交换图
\[
\begin{tikzcd}[row sep=large,column sep=large]
& a' \arrow[d] \arrow[ddl,"v_1"'] \arrow[ddr,"v_2"] & \\
& a'' \arrow[dl] \arrow[dr] & \\
a_1 \arrow[dr,"u_1"'] & & a_2 \arrow[dl,"u_2"] \\
& a &
\end{tikzcd}
\]
因此可将 \SGARefLink{sga7.sga7i.evi.n1-2}{1.2}(1) 替换为:

(1)$'$ 每个图
\[
\begin{tikzcd}[cramped,column sep=small,row sep=small]
a_1 \arrow[dr] & & a_2 \arrow[dl] \\
& a &
\end{tikzcd}
\qquad\text{在 }\underline A\text{ 中位于下述图之上}\qquad
\begin{tikzcd}[cramped,column sep=small,row sep=small]
R_1 \arrow[dr] & & R_2 \arrow[dl] \\
& R &
\end{tikzcd}
\quad\text{在 }\underline C_\Lambda,
\]
只要 $R_2\longrightarrow R$ 满射, 就能补为交换图
\[
\begin{tikzcd}[row sep=large,column sep=large]
& a' \arrow[dl] \arrow[dr] & \\
a_1 \arrow[dr] & & a_2 \arrow[dl] \\
& a &
\end{tikzcd}
\qquad\text{位于下述图之上}\qquad
\begin{tikzcd}[row sep=large,column sep=large]
& R_1\times_RR_2 \arrow[dl,"\pi_1"'] \arrow[dr,"\pi_2"] & \\
R_1 \arrow[dr] & & R_2 \arrow[dl] \\
& R &
\end{tikzcd}
\]
其中 $\pi_i$ ($i=1,2$) 为投影.

\underline{注 \SGARefTarget{sga7.sga7i.evi.n1-3.item-c}{1.3}(c).}
若 $\underline A$ 为 $\underline C_\Lambda$ 上的半齐性余纤维群胚,
则 $[\underline A(k[\varepsilon])]$ 典范地带有 $k$ 上的向量空间结构,
其中 $k[\varepsilon]$ 为 $k$ 上的对偶数环.
事实上, 设 $\underline{\underline V}$ 为 $k$ 上有限维向量空间所成的范畴.

%%%%%%%%%%  scan idx 47  |  folio 36  %%%%%%%%%%

每个保持积的函子 $H:\underline{\underline V}\longrightarrow(\mathrm{Ens})$
典范地分解为
\[
\begin{tikzcd}[row sep=large,column sep=large]
\underline{\underline V} \arrow[rr,"H"]
 \arrow[dr,"\overline H"'] & & (\mathrm{Ens}) \\
& \underline{\underline V} \arrow[ur] &
\end{tikzcd}
\]
其中 $\underline{\underline V}\longrightarrow(\mathrm{Ens})$ 为遗忘函子.
另一方面, 函子
\[
 G:\underline{\underline V}\longrightarrow\underline C_\Lambda,
 \qquad W\longmapsto k[W]=k\oplus W,\qquad W^2=0,
\]
把 $\underline{\underline V}$ 中的积送到 $\underline C_\Lambda$ 中在 $k$ 上的纤维积.
若 $\underline A$ 为 $\underline C_\Lambda$ 上的半齐性余纤维群胚,
则由 \SGARefLink{sga7.sga7i.evi.n1-3.item-a}{1.3}(a),
\[
 [\underline A(k[V]\times_k k[W])]
 \longrightarrow
 [\underline A(k[V])]\times[\underline A(k[W])]
\]
为双射. 因而函子 $V\longmapsto[\underline A(k[V])]$ 与积交换, 所证即得.

% END INCLUDED SGA7I_English_current_buildable_source_20260823/components/10_expose_VI_1_2_through_1_3.tex
% BEGIN INCLUDED SGA7I_English_current_buildable_source_20260823/components/11_expose_VI_1_4_through_1_6.tex
%%%%%%%%%%  scan idx 47  |  folio 36 (continued)  %%%%%%%%%%

\underline{定义 \SGARefTarget{sga7.sga7i.evi.n1-4}{1.4}.}
(1) \underline{设 $R$ 为 $\underline C_\Lambda$ 的对象.
$\underline C_\Lambda$ 中 $R$ 的一个小扩张(small extension)是满射
$R'\xrightarrow{f}R$ (在 $\underline C_\Lambda$ 中), 满足 $(\ker f)^2=0$ 且 $\ker f\simeq k$.
换言之, $R'\xrightarrow{f}R$ 为 $R$ 的小扩张, 若对某个满足 $t^2=0$ 的
$t\in R'$, 序列}
\[
0\longrightarrow k\xrightarrow{t}R'\xrightarrow{f}R
 \longrightarrow0
\]
\underline{正合.}

(2) \underline{设 $\underline A$ 为 $\underline C_\Lambda$ 上的余纤维群胚,
$a$ 为 $\underline A$ 的对象. $a$ 的小延拓(small prolongation)是
$\underline A$ 中满足 $P(\alpha)$ 为 $\underline C_\Lambda$ 中小扩张的箭头
$a'\xrightarrow{\alpha}a$. 称箭头 $a'\xrightarrow{\alpha}a$
对于 $a$ 的小延拓是遍有的(versal), 若对每个小延拓 $a'_1\longrightarrow a$,
存在交换图}
\[
\begin{tikzcd}[row sep=large,column sep=large]
a' \arrow[rr] \arrow[dr,"\alpha"'] & & a'_1 \arrow[dl] \\
& a &
\end{tikzcd}
\]

\underline{命题 \SGARefTarget{sga7.sga7i.evi.n1-5}{1.5}.
设 $\underline A$ 为 $\underline C_\Lambda$ 上的半齐性余纤维群胚.}

(1) \underline{设}
\[
\begin{tikzcd}[column sep=large]
& R \\
R\times_k k[\varepsilon]
 \arrow[ur,"\pi_1"] \arrow[dr,"\pi_2"'] & \\
& k[\varepsilon]
\end{tikzcd}
\]
\underline{为 $\underline C_\Lambda$ 中的投影映射,
$a'\xrightarrow{\alpha}a$ 为 $\pi_1$-映射.
则存在映射 $a\xrightarrow{\beta}a'$ 使 $\alpha\beta=1_a$,
当且仅当存在映射 $a\longrightarrow(\pi_2)_*a'$.}

%%%%%%%%%%  scan idx 48  |  folio 37  %%%%%%%%%%

(2) \underline{设 $R'\xrightarrow{f}R$ 为 $\underline C_\Lambda$ 中的小扩张,
并考虑交换图}
\[
\begin{tikzcd}[row sep=large,column sep=small]
& a'' \arrow[dl,"\alpha'"'] \arrow[dr,"\alpha'_1"]
  & & & & & R'\times_RR' \arrow[dl] \arrow[dr] & \\
a' \arrow[dr,"\alpha"'] & & a' \arrow[dl,"\alpha_1"]
  & & \underline{\text{位于右图之上}} & R' \arrow[dr] & & R' \arrow[dl] \\
& a & & & & & R &
\end{tikzcd}
\]
\underline{以}
\[
 \widetilde f(r'_1,r'_2)=r'_1(0)+(r'_1-r'_2).
\]
\underline{定义典范映射 $\widetilde f:R'\times_RR'\longrightarrow k[\ker f]$.
若存在映射 $a'\longrightarrow(\widetilde f)_*(a'')$,
则存在映射 $\gamma:a'\longrightarrow a'$ 使 $\alpha=\alpha_1\gamma$.}

\underline{证明.} (1) 必要性显然. 反之, 假定存在映射
$\varphi:a\longrightarrow(\pi_2)_*a'$. 置
\[
 h:R\longrightarrow R\times_k k[\varepsilon],
 \qquad h=1\times P(\varphi).
\]
复合 $\pi_1h$ 是恒等映射, 因此得到映射
\[
 a\mathrel{\mathop{\rightleftarrows}^{\gamma}_{\alpha_1}}h_*a,
 \qquad \alpha_1\gamma=1_a,
 \qquad P(\gamma)=h,
 \qquad P(\alpha_1)=\pi_1.
\]
另一方面, $\pi_2h=P(\varphi)$, 故存在 $\pi_2$-映射
\[
 h_*a\longrightarrow(\pi_2)_*a'.
\]
于是 $(\pi_2)_*(h_*a)=(\pi_2)_*a'$, 由 \SGARefLink{sga7.sga7i.evi.n1-2}{1.2}(2)
得到 $\underline C_\Lambda$-映射 $u:h_*a\longrightarrow a'$ 使 $\alpha u=\alpha_1$.
定义 $\beta=u\gamma$. 则
\[
 \alpha\beta=\alpha u\gamma=\alpha_1\gamma=1_a.
\]

(2) 映射
\[
 1\times\widetilde f:R'\times_RR'
 \longrightarrow R'\times_k k[\ker f]
\]
为同构, 与到第一个因子的投影相容, 且
$\pi_2(1\times\widetilde f)=\widetilde f$, 其中
$\pi_2:R'\times_k k[\ker f]\longrightarrow k[\ker f]$ 为第二投影.
因 $R'\xrightarrow fR$ 为小扩张, $k[\ker f]\simeq k[\varepsilon]$.
故若存在映射
\[
 a'\longrightarrow(\widetilde f)_*a''
   =(\pi_2)_*\bigl((1\times\widetilde f)_*a''\bigr),
\]
则由 (1) 得到映射 $\gamma':a'\longrightarrow a''$ 使 $\alpha'\gamma'=1_{a'}$.
定义 $\gamma=\alpha'_1\gamma':a'\longrightarrow a'$. 则 $\alpha_1\gamma=\alpha$.

%%%%%%%%%%  scan idx 49  |  folio 38  %%%%%%%%%%

\underline{推论 \SGARefTarget{sga7.sga7i.evi.n1-6}{1.6}.
设 $\underline A$ 为 $\underline C_\Lambda$ 上的半齐性余纤维群胚,
并考虑 $\underline A$ 中的图}
\[
\begin{tikzcd}[column sep=large]
a' \arrow[dr,"\alpha"'] & & a_1 \arrow[dl,"\alpha_1"] \\
& a &
\end{tikzcd}
\]
\underline{其中 $\alpha_1$ 是小延拓. 假定 $a$ 对于 $e$ 的小延拓是遍有的.
则存在映射 $\rho:a'\longrightarrow a_1$ 使 $\alpha=\alpha_1\rho$,
当且仅当 $\underline C_\Lambda$ 中存在映射
$f:P(a')\longrightarrow P(a_1)$ 使 $P(\alpha)=P(\alpha_1)f$.}

\underline{证明.} 假定 $P(\alpha)=P(\alpha_1)f$, 其中 $f:P(a')\longrightarrow P(a_1)$.
用 $a'$ 在 $f$ 下的正像替换它, 即可并且将假定 $P(\alpha)=P(\alpha_1)$.
置 $R'=P(a_1)$, $R=P(a)$. 因 $\underline A$ 半齐性, 得到交换图
\[
\begin{tikzcd}[row sep=large,column sep=small]
& a'' \arrow[dl,"\beta"'] \arrow[dr,"\beta_1"]
  & & & & & R'\times_RR' \arrow[dl] \arrow[dr] & \\
a' \arrow[dr,"\alpha"'] & & a_1 \arrow[dl,"\alpha_1"]
  & & \text{位于右图之上} & R' \arrow[dr] & & R' \arrow[dl] \\
& a & & & & & R &
\end{tikzcd}
\]
因 $a$ 对于 $e$ 的小延拓是遍有的, $a'$ 也如此;
于是由 \SGARefLink{sga7.sga7i.evi.n1-5}{1.5}(2) 得到所断言的结论.

% END INCLUDED SGA7I_English_current_buildable_source_20260823/components/11_expose_VI_1_4_through_1_6.tex

% BEGIN INCLUDED SGA7I_English_current_buildable_source_20260823/components/12_expose_VI_1_7_through_1_8.tex
%%%%%%%%%%  scan idx 49  |  folio 38 (continued)  %%%%%%%%%%

\underline{定义 \SGARefTarget{sga7.sga7i.evi.n1-7}{1.7}.
对 $\underline C_\Lambda$ 的每个对象 $R$, 置}
\[
 \overline\Omega_R=M/(m_\Lambda R+M^2),
\]
\underline{其中 $m_\Lambda$ 和 $M$ 分别为 $\Lambda$ 和 $R$ 的极大理想.
这是 $k$ 上的有限维向量空间, 且 $R\longmapsto\overline\Omega_R$ 定义函子}
\[
 \overline\Omega:\underline C_\Lambda
   \longrightarrow\underline{\underline V},
\]
\underline{其中 $\underline{\underline V}$ 为 $k$ 上有限维向量空间的范畴.
若 $\underline A$ 为 $\underline C_\Lambda$ 上的余纤维范畴,
则复合函子 $\underline A\xrightarrow P\underline C_\Lambda
\longrightarrow\underline{\underline V}$ 仍记为 $\overline\Omega$, 稍有记号滥用.}

\underline{命题 \SGARefTarget{sga7.sga7i.evi.n1-8}{1.8}.
设 $\underline A$ 为 $\underline C_\Lambda$ 上的半齐性余纤维群胚.
给定 $\underline A$ 中对于 $e$ 的小延拓遍有的对象 $a$,
存在箭头 $a'\longrightarrow a$ 满足:}

%%%%%%%%%%  scan idx 50  |  folio 39  %%%%%%%%%%

(i) \underline{$a'$ 对于 $a$ 的小延拓是遍有的;}

(ii) \underline{$P(a')\longrightarrow P(a)$ 满射,
其核被 $P(a)$ 的极大理想零化;}

(iii) \underline{$\overline\Omega_{a'}\longrightarrow\overline\Omega_a$ 为同构.}

\underline{证明.} 置 $R=P(a)$, 并考虑正合序列
\[
 0\longrightarrow J\longrightarrow S\longrightarrow R
 \longrightarrow0,
\]
其中 $S$ 为 $\Lambda$ 上有限个变量的形式幂级数环,
$J\subset M^2$, 而 $M$ 为 $S$ 的极大理想.
设 $\Sigma$ 为 $S$ 的满足下列条件的理想 $J_1$ 的集合:

(i) $J\supset J_1\supset MJ$, 且

(ii) $\underline A$ 中存在 $\pi_1$-态射 $a_1\longrightarrow a$,
其中 $\pi_1:S/J_1\longrightarrow R$ 为典范映射.

$\underline A$ 的性质 \SGARefLink{sga7.sga7i.evi.n1-2}{1.2}(1),
连同 $J/MJ$ 是 $k$ 上有限维向量空间这一事实, 蕴含 $\Sigma$ 有最小理想.
以 $J'$ 记之, 并取一个 $\pi'$-态射 $a'\xrightarrow g a$,
其中 $\pi':S/J'\longrightarrow R$ 为典范映射.
因 $J'\subset J\subset M^2$, 映射
$\overline\Omega_{a'}\longrightarrow\overline\Omega_a$ 为同构.

我们断言 $a'\xrightarrow g a$ 对于 $a$ 的小延拓是遍有的.
事实上, 设 $a_1\longrightarrow a$ 为小延拓, 置 $R_1=P(a_1)$.
得到行正合的交换图
\[
\begin{tikzcd}[column sep=small]
0 \arrow[r] & J/MJ \arrow[r] \arrow[d,"h"]
 & S/MJ \arrow[r] \arrow[d,"h_1"]
 & R \arrow[r] \arrow[d,equal] & 0 \\
0 \arrow[r] & k \arrow[r] & R_1 \arrow[r] & R \arrow[r] & 0.
\end{tikzcd}
\]
若 $h=0$, 则 $\pi_1:R_1\longrightarrow R$ 有截面,
因而 $h_1$ 诱导映射 $f:S/J'\longrightarrow R_1$.
若 $h\ne0$, 则 $h_1:S/MJ\longrightarrow R_1$ 满射,
故由 $J'$ 的定义, 仍诱导 $f:S/J'\longrightarrow R_1$.
两种情形下均得到交换图
\[
\begin{tikzcd}[column sep=large]
S/J' \arrow[rr,"f"] \arrow[dr,"\pi'"']
 & & R_1 \arrow[dl,"\pi_1"] \\
& R &
\end{tikzcd}
\]
因 $a$ 对于 $e$ 的小延拓是遍有的, 由推论 \SGARefLink{sga7.sga7i.evi.n1-6}{1.6}
得到交换图

%%%%%%%%%%  scan idx 51  |  folio 40  %%%%%%%%%%

\[
\begin{tikzcd}[column sep=large]
a' \arrow[rr,"\rho"] \arrow[dr] & & a_1 \arrow[dl] \\
& a &
\end{tikzcd}
\]
即所要求的图.

% END INCLUDED SGA7I_English_current_buildable_source_20260823/components/12_expose_VI_1_7_through_1_8.tex

% BEGIN INCLUDED SGA7I_English_current_buildable_source_20260823/components/13_expose_VI_pro_categories_through_1_10.tex
%%%%%%%%%%  scan idx 51  |  folio 40 (continued)  %%%%%%%%%%

设 $P:\underline A\longrightarrow\underline C_\Lambda$ 为以群胚为纤维的余纤维范畴.
它给出以群胚为纤维的余纤维范畴
\[
 \operatorname{pro}(P):\operatorname{pro}(\underline A)
 \longrightarrow\operatorname{pro}(\underline C_\Lambda)
\]
(前对象范畴 $\operatorname{pro}(\underline A)$ 的定义见
A. Grothendieck, S\'eminaire Bourbaki 195), 并纳入交换图
\[
\begin{tikzcd}[column sep=large,row sep=large]
\underline A \arrow[r] \arrow[d,"P"']
 & \operatorname{pro}(\underline A)
   \arrow[d,"\operatorname{pro}(P)"] \\
\underline C_\Lambda \arrow[r]
 & \operatorname{pro}(\underline C_\Lambda).
\end{tikzcd}
\]

设 $\widehat{\underline C}_\Lambda$ 为 $\operatorname{pro}(\underline C_\Lambda)$
的满子范畴, 其对象为系统 $(R_i)_{i\in\mathbb N}$, 其中
$R_i\in\underline C_\Lambda$, 每个映射 $\varphi_i:R_i\longrightarrow R_{i-1}$
满射, 且对所有充分大的 $i$, 它诱导同构
\[
 \frac{m_i}{m_\Lambda R_i+m_i^2}
 \xrightarrow{\ \sim\ }
 \frac{m_{i-1}}{m_\Lambda R_{i-1}+m_{i-1}^2},
\]
其中 $\mathbb N$ 表示自然数的有序范畴, $m_i$ 为 $R_i$ 的极大理想.

设 $\widehat{\underline A}$ 为 $\operatorname{pro}(\underline A)$ 的满子范畴,
由下式定义:
\[
 \operatorname{ob}(\widehat{\underline A})=
 \left\{a\in\operatorname{ob}(\operatorname{pro}(\underline A))
 \mathrel{\Big|}
 \operatorname{pro}(P)(a)\in
 \operatorname{ob}(\widehat{\underline C}_\Lambda)\right\}.
\]
得到以群胚为纤维的余纤维范畴
$\widehat P:\widehat{\underline A}\longrightarrow\widehat{\underline C}_\Lambda$
及交换图
\[
\begin{tikzcd}[column sep=large,row sep=large]
\underline A \arrow[r,hook] \arrow[d,"P"']
 & \widehat{\underline A} \arrow[d,"\widehat P"] \\
\underline C_\Lambda \arrow[r,hook]
 & \widehat{\underline C}_\Lambda.
\end{tikzcd}
\]
范畴 $\widehat{\underline C}_\Lambda$ 与剩余域固定为 $k$ 的、
在 $\Lambda$ 上形式有限型的完备局部 $\Lambda$-代数范畴典范等价;
此后将二者认同.

\underline{定义 \SGARefTarget{sga7.sga7i.evi.n1-9}{1.9}.
设 $\underline A$ 为 $\underline C_\Lambda$ 上的余纤维群胚.
称 $\widehat{\underline A}$ 的对象 $\xi$ 为拟始对象(quasi-initial),
若对 $\widehat{\underline A}$ 的每个对象 $a$, 存在 $\widehat{\underline A}$ 中的映射 $\xi\to a$.
称拟始对象 $\xi$ 为极小的, 若态射
$\xi:h_R\longrightarrow[\widehat{\underline A}]$ 诱导双射}

%%%%%%%%%%  scan idx 52  |  folio 41  %%%%%%%%%%

\[
 \xi(k[\varepsilon]):
 \operatorname{Hom}_{\Lambda\text{-alg}}(R,k[\varepsilon])
 \xrightarrow{\ \sim\ }
 [\underline A(k[\varepsilon])],
 \qquad R=P(\xi).
\]
\underline{称 $\widehat{\underline A}$ 的对象 $\xi$ 为 $\underline A$
的遍有形式对象(versal formal object), 若 $\underline A$ 中每个满射
$\alpha:a'\to a$ 诱导满射}
\[
 \operatorname{Hom}_{\widehat{\underline A}}(\xi,a')
 \longrightarrow
 \operatorname{Hom}_{\widehat{\underline A}}(\xi,a).
\]
\underline{称遍有形式对象 $\xi$ 为极小的, 若
$\xi:h_R\longrightarrow[\widehat{\underline A}]$ 在 $k[\varepsilon]$
上诱导上述双射, 其中 $R=P(\xi)$.}

\underline{命题 \SGARefTarget{sga7.sga7i.evi.n1-10}{1.10}.
设 $\underline A$ 为 $\underline C_\Lambda$ 上的余纤维群胚.}

(i) \underline{$\underline A$ 的每个(极小)遍有形式对象都是
$\widehat{\underline A}$ 的(极小)拟始对象.}

(ii) \underline{设 $\xi$ 为 $\widehat{\underline A}$ 的对象, 满足}
\[
 \widetilde\xi(k[\varepsilon]):
 \operatorname{Hom}_{\Lambda\text{-alg}}(R,k[\varepsilon])
 \longrightarrow[\underline A(k[\varepsilon])]
\]
\underline{为单射, 其中 $P(\xi)=R$. 则 $\xi$ 在 $\widehat{\underline A}$
中的每个自态射都是自同构. 特别地, $\widehat{\underline A}$ 的极小拟始对象
若存在, 则在非典范同构意义下唯一.}

\underline{证明.} (i) 设 $\xi$ 为 $\underline A$ 的遍有形式对象.
对 $\widehat{\underline A}$ 的对象 $a$, 取一个由 $\underline A$ 中满射给出的表示
\[
 \cdots\longrightarrow a_n\xrightarrow{\alpha_n}a_{n-1}
 \longrightarrow\cdots\longrightarrow
 a_1\xrightarrow{\alpha_1}a_0,
 \qquad a=\varprojlim_n a_n.
\]
假定已取定映射 $\varphi_i:\xi\longrightarrow a_i$, 对所有 $i<n$ 满足
$\alpha_i\varphi_i=\varphi_{i-1}$. 因 $\xi$ 遍有且 $\alpha_n$ 满射,
存在映射 $\varphi_n:\xi\longrightarrow a_n$ 使 $\alpha_n\varphi_n=\varphi_{n-1}$.
于是 $\varphi=\varprojlim_n\varphi_n$ 定义映射 $\xi\longrightarrow a$.

(ii) 设 $\alpha:\xi\longrightarrow\xi$ 为 $\widehat{\underline A}$ 中的自态射.
对每个 $\lambda\in\operatorname{Hom}_{\Lambda\text{-alg}}(R,k[\varepsilon])$, 有
\[
 \lambda_*\xi\simeq
 (\lambda\circ P(\alpha))_*\xi.
\]
因此由单射性, 对每个 $\lambda$ 有 $\lambda=\lambda\circ P(\alpha)$.
故 $P(\alpha):R\longrightarrow R$ 满射.
诺特 环的每个满射自态射都是自同构,
所以 $P(\alpha)$ 及 $\alpha$ 都是自同构.

% END INCLUDED SGA7I_English_current_buildable_source_20260823/components/13_expose_VI_pro_categories_through_1_10.tex
% BEGIN INCLUDED SGA7I_English_current_buildable_source_20260823/components/14_expose_VI_1_11_through_1_14.tex
%%%%%%%%%%  scan idx 52  |  folio 41 (continued)  %%%%%%%%%%

\underline{定理 \SGARefTarget{sga7.sga7i.evi.n1-11}{1.11} (Schlessinger).
设 $\underline A$ 为 $\underline C_\Lambda$ 上的余纤维群胚.
则 $\underline A$ 有极小遍有形式对象当且仅当:}

%%%%%%%%%%  scan idx 53  |  folio 42  %%%%%%%%%%

(i) \underline{$\underline A$ 半齐性;}

(ii) \underline{$[\underline A(k[\varepsilon])]$ 是 $k$ 上的有限维向量空间
(参见 \SGARefLink{sga7.sga7i.evi.n1-3.item-c}{1.3}(c)).}

\underline{证明. 必要性.}
设 $\xi\in\widehat{\underline A}$ 为 $\underline A$ 的极小遍有形式对象,
考虑 $\underline A$ 中的图
\[
\begin{tikzcd}[column sep=large,row sep=large]
a_1 \arrow[dr] & & a_2 \arrow[dl,"\alpha"] \\
& a &
\end{tikzcd}
\]
其中 $\alpha$ 满射. 因 $\xi$ 是拟始对象, 存在映射 $\xi\longrightarrow a_1$.
因 $\alpha$ 满射, 复合 $\xi\longrightarrow a_1\longrightarrow a$ 可提升到 $a_2$;
故得到交换图
\[
\begin{tikzcd}[column sep=large,row sep=large]
& \xi \arrow[dl] \arrow[dr] & \\
a_1 \arrow[dr] & & a_2 \arrow[dl,"\alpha"] \\
& a &
\end{tikzcd}
\]
这就证明了条件 \SGARefLink{sga7.sga7i.evi.n1-2}{1.2}(1).

现证明条件 \SGARefLink{sga7.sga7i.evi.n1-2}{1.2}(2). 设
\[
\begin{tikzcd}[column sep=large]
& S \\
S\times_k k[\varepsilon]
 \arrow[ur,"\pi_1"] \arrow[dr,"\pi_2"'] & \\
& k[\varepsilon]
\end{tikzcd}
\]
为投影, 并设
\[
\begin{tikzcd}[column sep=large]
a' \arrow[dr,"\alpha"'] & & b' \arrow[dl,"\beta"] \\
& a &
\end{tikzcd}
\]
为 $\underline A$ 中的 $\pi_1$-映射. 由 $\xi$ 的定义, 存在交换图
\[
\begin{tikzcd}[column sep=large,row sep=large]
& \xi \arrow[dl,"u"'] \arrow[dr,"v"] & \\
a' \arrow[dr,"\alpha"'] & & b' \arrow[dl,"\beta"] \\
& a &
\end{tikzcd}
\]
从而 $\pi_1P(u)=\pi_1P(v)$. 现假定 $(\pi_2)_*a'\simeq(\pi_2)_*b'$.
由 $\xi$ 的极小性得到 $\pi_2P(u)=\pi_2P(v)$, 因而

%%%%%%%%%%  scan idx 54  |  folio 43  %%%%%%%%%%

\[
 P(u)=\pi_1P(u)\times\pi_2P(u)
     =\pi_1P(v)\times\pi_2P(v)=P(v).
\]
因 $\underline A$ 为余纤维群胚, 存在 $\underline A(S\times_k k[\varepsilon])$
中唯一的同构 $\sigma:a'\xrightarrow{\sim}b'$ 使 $\sigma u=v$.
等式 $\beta\sigma u=\alpha u$ 于是蕴含 $\beta\sigma=\alpha$.

\underline{充分性.}
设 $c_1,\ldots,c_n$ 为 $\underline A(k[\varepsilon])$ 的对象,
其同构类构成 $[\underline A(k[\varepsilon])]$ 的一组 $k$-基.
取
\[
 a_1\in\underline A(
 k[\varepsilon]\times_k\cdots\times_k k[\varepsilon])
\]
使 $(\pi_j)_*a_1\simeq c_j$. 则
\[
 \widetilde a_1(k[\varepsilon]):
 \operatorname{Hom}_{\Lambda\text{-alg}}(R_1,k[\varepsilon])
 \longrightarrow[\underline A(k[\varepsilon])]
\]
为双射, 其中 $R_1=P(a_1)$. 特别地, $a_1\longrightarrow e$
对于 $e$ 的小延拓是遍有的. 取对于 $a_1$ 的小延拓遍有的
$a_2\longrightarrow a_1$, 使 $\overline\Omega_{a_2}\longrightarrow\overline\Omega_{a_1}$
为双射. 归纳继续, 取对于小延拓遍有且在 $\overline\Omega$ 上诱导双射的
$a_{n+1}\longrightarrow a_n$, 并置
\[
 \xi=\varprojlim_n a_n.
\]
则 $\overline\Omega_\xi\longrightarrow\overline\Omega_{a_1}$ 为双射, 故
\[
 \xi(k[\varepsilon]):
 \operatorname{Hom}_{\Lambda\text{-alg}}(R,k[\varepsilon])
 \xrightarrow{\ \sim\ }
 [\underline A(k[\varepsilon])]
\]
为双射.

还须证明遍有性. 考虑图
\[
\begin{tikzcd}
& a' \arrow[d,"\alpha"] \\
\xi \arrow[r,"h"] & a
\end{tikzcd}
\]
其中 $\alpha$ 满射. 须将 $h$ 提升到 $a'$.
用归纳法可假定 $a'$ 为 $a$ 的小延拓. 因 $P(a)$ 是 Artin 环,
$h$ 经过某个 $a_q$ 分解,
\[
 \xi\longrightarrow a_q\longrightarrow a.
\]
半齐性给出交换图
\[
\begin{tikzcd}
& a'' \arrow[r] \arrow[d,"\alpha'"] & a' \arrow[d,"\alpha"] \\
\xi \arrow[r] & a_q \arrow[r] & a,
\end{tikzcd}
\]
其中 $\alpha'$ 也为小延拓. 因 $a_{q+1}\longrightarrow a_q$
对于 $a_q$ 的小延拓遍有, 存在交换图
\[
\begin{tikzcd}
& & a'' \arrow[d,"\alpha'"] \\
\xi \arrow[r] & a_{q+1} \arrow[r] \arrow[ur] & a_q,
\end{tikzcd}
\]
它给出所需的提升.

\underline{注 \SGARefTarget{sga7.sga7i.evi.n1-12}{1.12}.} 设

%%%%%%%%%%  scan idx 55  |  folio 44  %%%%%%%%%%

\[
 \cdots\longrightarrow a_n\longrightarrow a_{n-1}
 \longrightarrow\cdots\longrightarrow a_2\longrightarrow a_1
 \longrightarrow a_0=e
\]
为 $\underline A$ 中的无穷箭头序列, 满足
\[
\left\{
\begin{aligned}
&a_i\text{ 对于小延拓遍有, 对象为 }a_{i-1},\\
&P(a_i)\longrightarrow P(a_{i-1})\text{ 满射},\\
&\overline\Omega_{a_i}\longrightarrow
 \overline\Omega_{a_{i-1}}\text{ 为双射, 当 }i\gg0.
\end{aligned}
\right.
\]
上述证明表明 $\varprojlim_i a_i$ 是 $\underline A$ 的遍有形式对象.

\underline{注 \SGARefTarget{sga7.sga7i.evi.n1-13}{1.13}.}
设 $\underline A$ 为 $\underline C_\Lambda$ 上的余纤维群胚.
若 $\underline A$ 满足 \SGARefLink{sga7.sga7i.evi.n1-11}{1.11}(i) 及 (ii),
则 $[\underline A]$ 亦然; 而且由 1.10(ii),
$\xi\in\operatorname{ob}(\widehat{\underline A})$ 为 $\underline A$
的极小遍有形式对象, 当且仅当它对 $[\underline A]$ 亦然.

\underline{注 \SGARefTarget{sga7.sga7i.evi.n1-14}{1.14}.}
设 $\underline A$ 为 $\underline C_\Lambda$ 上有极小遍有形式对象 $\xi$
的余纤维群胚. $R=P(\xi)$ 是在 $\Lambda$ 上形式有限型的完备局部
$\Lambda$-代数, 它的每个不变量自动都是 $\underline A$
在 $\underline C_\Lambda$ 上的不变量.
例如, 若 $R$ 在 $\Lambda$ 上(形式)光滑, 即为 $\Lambda$ 上的形式幂级数代数,
则称 $\underline A$ 在 $\underline C_\Lambda$ 上(形式)光滑;
若 $\operatorname{Spec}(R)$ 不可约, 则称 $\underline A$
在 $\underline C_\Lambda$ 上不可约.

下一节将对 $\underline C_\Lambda$ 上的函子
$\underline A\longrightarrow\underline B$ 引入光滑性,
其中 $\underline A$ 和 $\underline B$ 是 $\underline C_\Lambda$
上以群胚为纤维的余纤维范畴.
由 \SGARefLink{sga7.sga7i.evi.n2-4}{2.4}, $\underline A$
在 $\underline C_\Lambda$ 上(形式)光滑等价于函子
$P:\underline A\longrightarrow\underline C_\Lambda$ 光滑, 故不会发生歧义.
定义
\[
 \dim_{\underline C_\Lambda}\underline A
   =\dim(k\otimes_\Lambda R).
\]

为此, 设 $\underline C_k$ 为剩余域为 $k$ 的 Artin 局部 $k$-代数范畴.
函子 $\underline C_k\longrightarrow\underline C_\Lambda$ 为全忠实函子,
其左伴随为
\[
 \underline C_\Lambda\longrightarrow\underline C_k,
 \qquad R\longmapsto k\otimes_\Lambda R.
\]
设 $\underline A|_{\underline C_k}$ 为
$P:\underline A\longrightarrow\underline C_\Lambda$
沿 $\underline C_k\longrightarrow\underline C_\Lambda$ 的拉回.
若 $\xi$ 为 $\underline A$ 在 $\underline C_\Lambda$ 上的极小遍有形式对象,
则 $k\otimes_\Lambda\xi$ 对于 $\underline A|_{\underline C_k}$
在 $\underline C_k$ 上极小遍有;
箭头 $\xi\longrightarrow k\otimes_\Lambda\xi$
在 $R\longrightarrow k\otimes_\Lambda R$ 上余笛卡尔. 从而
\[
 \dim_{\underline C_\Lambda}\underline A
 =\dim_{\underline C_k}(\underline A|_{\underline C_k}).
\]

% END INCLUDED SGA7I_English_current_buildable_source_20260823/components/14_expose_VI_1_11_through_1_14.tex

% BEGIN INCLUDED SGA7I_English_current_buildable_source_20260823/components/15_expose_VI_1_15_through_1_17.tex
%%%%%%%%%%  scan idx 56  |  folio 45  %%%%%%%%%%

\textbf{\SGARefTarget{sga7.sga7i.evi.n1-15}{1.15}.}
本段余下部分简要说明如何修改上述理论以容许剩余域扩张.

设 $\Lambda$ 和 $K$ 为完备诺特局部环, $K$ 为 $\Lambda$-代数. 假定:

(a) 自然映射 $\Lambda\longrightarrow K$ 为局部同态;

(b) $K$ 的剩余域 $K'$ 为 $\Lambda$ 的剩余域 $k$ 的有限型扩张.

最有意义的情形是 $K=K'$, 且 $K$ 是 $k$ 的有限不可分扩张.

以 $\widehat C_{\Lambda,K}$ 表示如下范畴: 对象为配有
$\Lambda$-代数结构及 $\Lambda$-满同态 $s:R\longrightarrow K$
的完备诺特局部环 $R$. 设 $C_{\Lambda,K}$ 为其满子范畴,
由满足 $\ker(s)$ 作为 $R$-模长度有限的二元组 $(R,s)$ 组成.
范畴 $\widehat C_{\Lambda,K}$ 可认同为
$\operatorname{Pro}(C_{\Lambda,K})$ 的满子范畴.
若 $K$ 为 Artin 环, 则 $C_{\Lambda,K}$ 的每个对象都是 Artin 环.

对 $K'$ 上的有限维向量空间 $V$, 用 $D_V(K)$
表示由 $V$ 定义的 $K$ 上的对偶数代数.
其元素形如 $x+v\varepsilon$, 其中 $x\in K$, $v\in V$, $\varepsilon^2=0$,
乘法为
\[
 (x+v\varepsilon)(x'+v'\varepsilon)
 =xx'+(xv'+x'v)\varepsilon.
\]
映射
\[
 s:D_V(K)\longrightarrow K,
 \qquad s(x+v\varepsilon)=x,
\]
使 $D_V(K)$ 成为 $C_{\Lambda,K}$ 的对象.
此前 $k[\varepsilon]$ 的作用现由 $D_{K'}(K)$ 承担.

如 \SGARefLink{sga7.sga7i.evi.n1-8}{1.8}, $C_{\Lambda,K}$ 上的每个余纤维群胚
$\underline A$ 自然延伸为 $\widehat C_{\Lambda,K}$ 上的余纤维范畴
$\widehat{\underline A}$. 若 $\underline A$ 或 $\widehat{\underline A}$
中的映射所对应的 $C_{\Lambda,K}$ 或 $\widehat C_{\Lambda,K}$
中的代数同态满射, 就称该映射\underline{满射}.

%%%%%%%%%%  scan idx 57  |  folio 46  %%%%%%%%%%

\underline{假定 $\underline A(K)=\{e\}$.}

\underline{定义 \SGARefTarget{sga7.sga7i.evi.n1-16}{1.16}.
称 $C_{\Lambda,K}$ 上的余纤维群胚 $\underline A$ 为半齐性的, 若满足下列性质:}

(I) \underline{$\underline A$ 中每个由实线箭头组成的图}
\[
\begin{tikzcd}[row sep=large,column sep=large]
a'' \arrow[r,dashed] \arrow[d,dashed]
 & a' \arrow[d] \\
a_1 \arrow[r] & a
\end{tikzcd}
\]
\underline{只要 $a'\longrightarrow a$ 满射, 就能添加点线箭头补为交换图.}

(II) \underline{对每个 $R$, 自然映射}
\[
 [\underline A(R\times_KD_{K'}(K))]
 \longrightarrow
 [\underline A(R)]\times[\underline A(D_{K'}(K))]
\]
\underline{为双射.}

若 $\underline A$ 半齐性, 则 $[\underline A]$ 也半齐性.

\textbf{\SGARefTarget{sga7.sga7i.evi.n1-17}{1.17}.}
先将 $\underline A$ 限制到 $C_{\Lambda,K}$ 中由 $D_V(K)$ 组成的满子范畴.
置
\[
 \Omega=\Omega^1_{K/\Lambda}\otimes_KK',
 \qquad d:K\longrightarrow\Omega.
\]
映射 $f:D_V(K)\longrightarrow D_W(K)$ 可写成
\[
 f(x+v\varepsilon)
 =x+\bigl(u(v)+D(dx)\bigr)\varepsilon,
\]
其中 $u\in\operatorname{Hom}_{K'}(V,W)$,
$D\in\operatorname{Hom}_{K'}(\Omega,W)$.
将 $f$ 与 $(u,D)$ 认同, 则复合由下式给出:
\[
 (u,D)\circ(u',D')=(uu',D+uD').
\]
条件 (II) 蕴含:

(II$'$) 从 $K'$-向量空间到集合的函子
\[
 V\longmapsto[\underline A(D_V(K))]
\]
与积交换.

% END INCLUDED SGA7I_English_current_buildable_source_20260823/components/15_expose_VI_1_15_through_1_17.tex

% BEGIN INCLUDED SGA7I_English_current_buildable_source_20260823/components/16_expose_VI_lemma_1_18.tex
% R3 explicit revisions VI-A-01; see qa/FIDELITY-R3-CHANGES.json.
%%%%%%%%%%  scan idx 58  |  folio 47  %%%%%%%%%%

\underline{引理 \SGARefTarget{sga7.sga7i.evi.n1-18}{1.18}.
假定 $\underline A$ 满足条件 (II$'$). 则:}

(1) $[\underline A(D_{K'}(K))]$ 带有唯一的向量空间结构,
使得关于 $V$ 函子性地有
\[
 [\underline A(D_V(K))]
 \simeq
 V\otimes_{K'}[\underline A(D_{K'}(K))].
\]

(2) 存在线性映射
\[
 \gamma:\operatorname{Hom}(\Omega,K')
 \longrightarrow [\underline A(D_{K'}(K))]
\]
使对每个 $D\in\operatorname{Hom}(\Omega,K')$, 映射
\[
 [\underline A((1,D))]:
 [\underline A(D_{K'}(K))]
 \longrightarrow [\underline A(D_{K'}(K))]
\]
为 $x\longmapsto x+\gamma(D)$.

(3) 对每个 $(u,D):D_V(K)\longrightarrow D_W(K)$, 诱导映射
\[
 (u,D):[\underline A(D_V(K))]
 \longrightarrow[\underline A(D_W(K))],
\]
即
\[
 (u,D):
 V\otimes[\underline A(D_{K'}(K))]
 \longrightarrow
 W\otimes[\underline A(D_{K'}(K))],
\]
由
\[
 x\longmapsto u(x)+\gamma(D).
\]
给出.

引入记号
\[
\begin{aligned}
 F(V)&=[\underline A(D_V(K))],\\
 A&=F(K'),\\
 F(u,D):F(V)&\longrightarrow F(W)
   &&\text{表示下述映射诱导的映射: }
     (u,D):D_V(K)\longrightarrow D_W(K),\\
 F(D):F(V)&\longrightarrow F(V),
   &F(D)&=F(1,D).
\end{aligned}
\]
由 (II$'$), 并因 $K'$ 是 $K'$-向量空间范畴中的 $K'$-向量空间对象,
$F(K')$ 是集合范畴中的 $K'$-向量空间对象, 故为 $K'$-向量空间.
这证明了 (1). 此后将 $F(V)$ 与 $V\otimes A$ 认同.

恒等式
\[
 (u,D)=(1,D)\circ(u,0)
\]
给出
\[
 F(u,D)=F(D)\circ(u\otimes1_A).
\]
$F$ 的函子性等价于恒等式
\[
\begin{aligned}
 F(0)&=\operatorname{Id},\\
 F(D)\circ F(D')&=F(D+D'),\\
 (u\otimes1)\circ F(D)
   &=F(uD)\circ(u\otimes1_A).
\end{aligned}
\]

对任意 $D:\Omega\longrightarrow V$ 及任意 $n\in\mathbb N$, 设
\[
 i:V\lhook\joinrel\longrightarrow V\oplus K'^{\,n},
 \qquad
 p:V\oplus K'^{\,n}\longrightarrow K'^{\,n}
\]
为典范的嵌入及投影. 则 $piD=0$, 且下图交换:

%%%%%%%%%%  scan idx 59  |  folio 48  %%%%%%%%%%

\[
\begin{tikzcd}[column sep=huge,row sep=large]
 V\otimes A
   \arrow[r]
   \arrow[d,"{F(D)}"]
 & (V\oplus K'^{\,n})\otimes A
   \arrow[r]
   \arrow[d,"{F(i,D)}"]
 & K'^{\,n}\otimes A
   \arrow[d,equal]\\
 V\otimes A
   \arrow[r]
 & (V\oplus K'^{\,n})\otimes A
   \arrow[r]
 & K'^{\,n}\otimes A.
\end{tikzcd}
\]

设 $x:K'^{\,n}\longrightarrow V$ 为映射, 置 $x_i=x(e_i)$,
并设 $(a_i)$ 为 $A$ 中的一族元素. 有
\[
 (1_V+x)\circ i\circ D=D.
\]
从而, 若
\[
 F(iD)\Bigl(\sum_i e_i a_i\Bigr)
 =\sum_i e_i a_i+a,
 \qquad a\in V\otimes A,
\]
则
\[
 F(D)\Bigl(\sum_i x_i a_i\Bigr)
 =\sum_i x_i a_i+a,
\]
且 $a$ 与 $x_i$ 无关. 故
\[
\begin{aligned}
 F(D)(y)&=y+G(D),\\
 G(D+D')&=G(D)+G(D'),\\
 uG(D)&=G(uD)
   \qquad\text{对每个 }u:V\longrightarrow W.
\end{aligned}
\]
每个这样的 $G$ 均由某个 $\gamma\in\Omega\otimes A$ 按下式给出:
\[
 G(D)=(D\otimes1_A)(\gamma).
\]
这就证明了 (2) 及 (3).

% END INCLUDED SGA7I_English_current_buildable_source_20260823/components/16_expose_VI_lemma_1_18.tex
% BEGIN INCLUDED SGA7I_English_current_buildable_source_20260823/components/17_expose_VI_1_19_through_1_20.tex
%%%%%%%%%%  scan idx 59  |  folio 48 (continued)  %%%%%%%%%%

\underline{定义 \SGARefTarget{sga7.sga7i.evi.n1-19}{1.19}.}
(1) 称对象 $\xi\in\widehat{\underline A}$ 为\underline{遍有形式对象},
若 $\underline A$ 中的每个满射 $\alpha:a'\longrightarrow a$ 诱导满射
\[
 \operatorname{Hom}_{\widehat{\underline A}}(\xi,a')
 \longrightarrow
 \operatorname{Hom}_{\widehat{\underline A}}(\xi,a).
\]

(2) 对于在 $\widehat C_{\Lambda,K}$ 中的像为 $(R,s)$ 的遍有形式对象
$\xi\in\widehat{\underline A}$, 称其为\underline{极小的}, 若满足下列条件.
对任意两个映射
\[
 \varphi_1,\varphi_2:R\rightrightarrows D_{K'}(K),
\]
若 $\varphi_1(\xi)$ 与 $\varphi_2(\xi)$ 同构,
则存在导子 $D:K\longrightarrow K'$ 使
\[
 (1,D)(\varphi_1)=\varphi_2,
 \qquad\text{即}\qquad
 \varphi_2-\varphi_1=Ds.
\]

\underline{定理 \SGARefTarget{sga7.sga7i.evi.n1-20}{1.20}.
若 $\underline A$ 半齐性, 且}
\[
 \dim_{K'}[\underline A(D_{K'}(K))]<\infty,
\]
\underline{则 $\underline A$ 有极小遍有形式对象, 且任意两个这样的对象都同构.}

%%%%%%%%%%  scan idx 60  |  folio 49  %%%%%%%%%%

沿用 \SGARefLink{sga7.sga7i.evi.n1-18}{1.18} 的记号,
设 $H^*$ 为 $[\underline A(D_{K'}(K))]$ 中 $\gamma$ 的像的补空间,
$H$ 为 $H^*$ 的对偶. 置
\[
 R_0=D_H(K).
\]
取 $\xi_0\in\underline A(R_0)$ 使其同构类在认同
\[
 [\underline A(R_0)]
 \simeq H\otimes[\underline A(D_{K'}(K))]
 \simeq
 \operatorname{Hom}\bigl(H^*,[\underline A(D_{K'}(K))]\bigr),
\]
下为包含映射 $H^*\hookrightarrow[\underline A(D_{K'}(K))]$. 则:

(1) 对每个 $V$, 映射
\[
 \operatorname{Hom}_{C_{\Lambda,K}}
   \bigl(D_H(K),D_V(K)\bigr)
 \longrightarrow[\underline A(D_V(K))]
\]
满射;

(2) $(R_0,\xi_0)$ 满足 \SGARefLink{sga7.sga7i.evi.n1-19}{1.19}(2) 的极小性条件.

取在 $\Lambda$ 上形式光滑的 $S\in\widehat C_{\Lambda,K}$,
并取满同态 $t:S\longrightarrow R_0$, 使
\[
 \Omega_{R_0/\Lambda}\otimes K'
 \xrightarrow{\ \sim\ }
 \Omega_{S/\Lambda}\otimes K'.
\]
若 $M$ 为 $S$ 的极大理想, 则对 $n\geq1$, 用下列条件归纳定义
\[
 R_n=S/I_n,
 \qquad
 \xi_n\in\operatorname{Ob}(\underline A(R_n)),
\]

(a) $I_{n+1}$ 为所有满足
\[
 I_nM\subset I\subset I_n
\]
且 $\xi_n$ 可延伸为对象 $\xi\in\underline A(S/I)$ 的理想 $I$ 的交;

(b) $\xi_{n+1}$ 为 $\xi_n$ 的延伸.

若
\[
 R=\varprojlim_nR_n,
 \qquad
 \xi=\varprojlim_n\xi_n,
\]
则 $(R,\xi)$ 为极小遍有形式对象, 并在同构意义下唯一.

% END INCLUDED SGA7I_English_current_buildable_source_20260823/components/17_expose_VI_1_19_through_1_20.tex

% BEGIN INCLUDED SGA7I_English_current_buildable_source_20260823/components/18_expose_VI_section_2_opening_through_2_2.tex
\section*{\SGARefTarget{sga7.sga7i.evi.n2}{2}. 前可表的余纤维群胚}

设 $\underline A$ 为 $\underline C_\Lambda$ 上的余纤维群胚,
有极小遍有形式对象 $\xi$. 一般而言, 单凭 $\xi$
不能在 $\underline C_\Lambda$-等价意义下确定 $\underline A$.
还须知道两个映射
\[
 R\mathrel{\substack{\xrightarrow{\ f\ }\\[-.6ex]
                      \xrightarrow[\ g\ ]{}}}S
\]
何时给出 $\underline C_\Lambda$-同构
$\xi(f)\longrightarrow\xi(g)$, 其中 $R=P(\xi)$.
作为 \SGARefLink{sga7.sga7i.evi.n1-11}{1.11} 及
\SGARefLink{sga7.sga7i.evi.n1-12}{1.12} 的应用,
现考虑在 $\underline C_\Lambda$-等价意义下恢复 $\underline A$ 的问题.
对一大类余纤维群胚, 将得到令人满意的答案.

%%%%%%%%%%  scan idx 61  |  folio 50  %%%%%%%%%%

实际出现的许多余纤维群胚满足一个略强于半齐性的性质;
将在 \SGARefLink{sga7.sga7i.evi.n2-5}{2.5} 明确定义它.
为研究这些群胚, 用“光滑函子”来重新表述上面的讨论是方便的.
本节继续使用上一节的所有记号与术语.

\underline{定义 \SGARefTarget{sga7.sga7i.evi.n2-1}{2.1}.}
设 $\varphi:\underline A\longrightarrow\underline B$
为 $\underline C_\Lambda$ 上余纤维群胚间的函子.
若 $\underline B$ 中每个满射 $\beta:b'\longrightarrow\varphi(a)$
都能补成交换图
\begin{equation}\tag{*}
\begin{tikzcd}[row sep=large,column sep=huge]
 \varphi(a') \arrow[r] \arrow[dr,"{\varphi(\alpha)}"']
 & b' \arrow[d,"\beta"]\\
 & \varphi(a)
\end{tikzcd}
\end{equation}
其中 $\alpha:a'\longrightarrow a$ 为 $\underline A$ 中某个映射,
则称 $\varphi$ 为\underline{光滑的}.
若 $\varphi$ 光滑且
\[
 \varphi(k[\varepsilon]):
 [\underline A(k[\varepsilon])]
 \longrightarrow
 [\underline B(k[\varepsilon])]
\]
为双射, 则称其为\underline{极小光滑的}.

\underline{注 \SGARefTarget{sga7.sga7i.evi.n2-2.item-a}{\SGARefAliasTarget{sga7.sga7i.evi.n2-2}2.2}(a).}
设 $\varphi:\underline A\longrightarrow\underline B$
为 $\underline C_\Lambda$ 上余纤维群胚间的函子. 下列条件等价:

(i) $\varphi$ 光滑;

(ii) $\underline B$ 中每个满射 $\beta:b'\longrightarrow\varphi(a)$
都能补成图 $(*)$, 其中 $\varphi(a')\longrightarrow b'$
为 $\underline C_\Lambda$-同构;

(iii) $\underline B$ 中每个小延拓 $\beta:b'\longrightarrow\varphi(a)$
都能补成交换图 $(*)$.

\underline{注 \SGARefTarget{sga7.sga7i.evi.n2-2.item-b}{2.2}(b).}
对 $\underline C_\Lambda$ 上每个余纤维群胚 $\underline A$, 自然函子
\[
 \underline A\longrightarrow[\underline A]
\]
极小光滑.

\underline{注 \SGARefTarget{sga7.sga7i.evi.n2-2.item-c}{2.2}(c).}
对每个 $R\in\widehat{\underline C}_\Lambda$, 设
\[
 h_R:\underline C_\Lambda\longrightarrow(\mathrm{Ens})
\]
为由
\[
 h_R(S)=\operatorname{Hom}_{\Lambda\text{-alg}}(R,S)
 \qquad(S\in\underline C_\Lambda).
\]
定义的函子. 它确定 $\underline C_\Lambda$ 上的余纤维离散群胚 $\underline h_R$.
$\widehat{\underline C}_\Lambda$ 中的映射 $R\longrightarrow S$ 诱导函子
\[
 \underline h_S\longrightarrow\underline h_R.
\]
此函子光滑当且仅当 $S$ 是 $R$ 上的形式幂级数环.

%%%%%%%%%%  scan idx 62  |  folio 51  %%%%%%%%%%

\underline{注 \SGARefTarget{sga7.sga7i.evi.n2-2.item-d}{2.2}(d).}
给定 $R\in\widehat{\underline C}_\Lambda$, 用 $(R,\underline C_\Lambda)$
表示 $\widehat{\underline C}_\Lambda$ 中源为 $R$、靶在
$\underline C_\Lambda$ 内的箭头范畴.
设 $\underline A$ 为 $\underline C_\Lambda$ 上的余纤维群胚.
给定 $P(\xi)=R$ 的 $\xi\in\operatorname{Ob}\widehat{\underline A}$, 设
\[
 \underline A'\longrightarrow(R,\underline C_\Lambda)
\]
为 $\underline A\longrightarrow\underline C_\Lambda$
沿典范函子 $(R,\underline C_\Lambda)\longrightarrow\underline C_\Lambda$ 的拉回.
取截面 $\xi^\#$ 使 $\xi^\#(1_R)=\xi$.
即对每个映射 $f:R\longrightarrow S$, 选取一个 $f$-映射
$\xi\longrightarrow\xi(f)$.

将 $(R\xrightarrow{f}S)$ 送至 $\xi(f)$, 得到 $\underline C_\Lambda$ 上的函子
\[
 \xi^\#:\underline h_R\longrightarrow\underline A
\]
在自然等价意义下, 此函子与 $\xi^\#$ 的选择无关.

断言 $\xi^\#:\underline h_R\longrightarrow\underline A$
(极小)光滑, 就是说 $\underline A$ 中的每个满射
$a'\xrightarrow{\alpha}\xi(f)$ 都能补成交换图
\[
\begin{tikzcd}[row sep=large,column sep=huge]
 \xi(f') \arrow[r] \arrow[dr,"{\xi(\,)}"']
 & a' \arrow[d,"\alpha"]\\
 & \xi(f)
\end{tikzcd}
\]
(并且在极小情形下,
$\underline h_R(k[\varepsilon])\longrightarrow\underline A(k[\varepsilon])$ 为双射).
这恰好就是说 $\xi$ 为 $\underline A$ 在 $\underline C_\Lambda$
上的(极小)遍有形式对象.

因此, $\underline A$ 有(极小)遍有形式对象等价于对某个
$R\in\operatorname{Ob}(\widehat{\underline C}_\Lambda)$,
存在 $\underline C_\Lambda$ 上的(极小)光滑函子
\[
 \underline h_R\longrightarrow\underline A
\]

% END INCLUDED SGA7I_English_current_buildable_source_20260823/components/18_expose_VI_section_2_opening_through_2_2.tex

% BEGIN INCLUDED SGA7I_English_current_buildable_source_20260823/components/19_expose_VI_2_3_through_2_7.tex
% R3 explicit revisions VI-A-04; see qa/FIDELITY-R3-CHANGES.json.
下述光滑性的一般性质是直接的.

\underline{命题 \SGARefTarget{sga7.sga7i.evi.n2-3}{2.3}.} 设
\[
 \underline A\xrightarrow{\ \varphi\ }\underline B
 \xrightarrow{\ \psi\ }\underline C
\]
为 $\underline C_\Lambda$ 上余纤维群胚间的函子.

(a) 若 $\varphi$ 和 $\psi$ 均光滑, 则 $\psi\circ\varphi$ 也光滑.

(b) 若 $\psi\circ\varphi$ 光滑且 $\varphi$ 在对象同构类上满射, 则 $\psi$ 光滑.

%%%%%%%%%%  scan idx 63  |  folio 52  %%%%%%%%%%

\underline{注 \SGARefTarget{sga7.sga7i.evi.n2-4}{2.4}.}
设 $\underline A$ 为 $\underline C_\Lambda$ 上有遍有形式对象 $\xi$ 的余纤维群胚.
由 \SGARefLink{sga7.sga7i.evi.n2-2.item-d}{2.2}(d), 这就是说
\[
 \xi:\underline h_R\longrightarrow\underline A
\]
光滑, 其中 $R=P(\xi)$. 因而下列断言等价:

\begin{itemize}
\item $\underline A$ 在 $\underline C_\Lambda$ 上形式光滑
      (参见 \SGARefLink{sga7.sga7i.evi.n1-14}{1.14}),
      即 $R$ 是 $\Lambda$ 上的形式幂级数代数;
\item 复合函子
      \[
       \underline h_R\xrightarrow{\ \xi\ }\underline A
       \xrightarrow{\ P\ }\underline C_\Lambda
       \;(=\underline h_\Lambda)
      \]
      光滑;
\item $P:\underline A\longrightarrow\underline C_\Lambda$ 光滑.
\end{itemize}
最后一个等价性由命题 \SGARefLink{sga7.sga7i.evi.n2-3}{2.3}(b) 得到.

现在考虑一个强于半齐性的性质.

\underline{定义 \SGARefTarget{sga7.sga7i.evi.n2-5}{2.5}.}
称 $\underline C_\Lambda$ 上的余纤维群胚 $\underline A$ 为\underline{齐性的},
若每个图
\[
\begin{tikzcd}[row sep=large,column sep=huge]
 & a' \arrow[d]\\
 b' \arrow[r] & a
\end{tikzcd}
\]
只要 $a'\longrightarrow a$ 满射, 就有 $\underline A$ 中的纤维积
$a'\times_a b'$, 且
\[
 P(a'\times_a b')=
 P(a')\times_{P(a)}P(b').
\]

\underline{注 \SGARefTarget{sga7.sga7i.evi.n2-6.item-a}{\SGARefAliasTarget{sga7.sga7i.evi.n2-6}2.6}(a).}
设 $\underline A$ 为 $\underline C_\Lambda$ 上的余纤维群胚, 考虑其中的图
\[
(*)\qquad
\begin{tikzcd}[row sep=small,column sep=small]
 & a' \arrow[dl] \arrow[dr] & \\
 a_1 \arrow[dr] & & a_2 \arrow[dl]\\
 & a &
\end{tikzcd}
\qquad\text{位于下述图之上}\qquad
\begin{tikzcd}[row sep=small,column sep=small]
 & R_1\times_R R_2 \arrow[dl,"{\pi_1}"'] \arrow[dr,"{\pi_2}"] & \\
 R_1 \arrow[dr] & & R_2 \arrow[dl]\\
 & R &
\end{tikzcd}
\]
右图在 $\underline C_\Lambda$ 中, $\pi_i$ 为投影.
若 $\underline A$ 中存在纤维积 $a_1\times_a a_2$, 则存在与投影相容的唯一映射
\[
 \varphi:a'\longrightarrow a_1\times_a a_2
\]
若
\[
 P(a_1\times_a a_2)=R_1\times_R R_2,
\]
则 $P(\varphi)$ 是 $R_1\times_R R_2$ 的恒等映射;
故 $\varphi$ 是 $\underline C_\Lambda$-同构, 即 $a'=a_1\times_a a_2$.
因此, 若 $\underline A$ 齐性, 则只要 $R_1\longrightarrow R$ 满射,
每个图 $(*)$ 均定义纤维积 $a_1\times_a a_2$.
立即得到

%%%%%%%%%%  scan idx 64  |  folio 53  %%%%%%%%%%

每个齐性余纤维群胚必然半齐性.

\underline{注 \SGARefTarget{sga7.sga7i.evi.n2-6.item-b}{2.6}(b).}
考虑范畴及函子所成的图
\[
\begin{tikzcd}[row sep=large,column sep=large]
 X \arrow[dr,"{\varphi}"'] & & Y \arrow[dl,"{\psi}"]\\
 & Z &
\end{tikzcd}
\]
通常的纤维积 $X\times_ZY$ 为如下范畴:
对象为满足 $\varphi(x)=\psi(y)$ 的二元组 $(x,y)$,
映射 $(x,y)\longrightarrow(x',y')$ 为满足
$\varphi(\alpha)=\psi(\beta)$ 的二元组
\[
 \alpha:x\longrightarrow x',\qquad
 \beta:y\longrightarrow y'
\]
$2$-范畴意义下的纤维积 $X\times_Z^2Y$ 为如下范畴:
对象为三元组 $(x,y;\gamma)$, 其中
\[
 \varphi(x)\xrightarrow{\ \gamma\ }\psi(y)
\]
为同构. 映射 $(x,y;\gamma)\longrightarrow(x',y';\gamma')$
为使下图交换的二元组 $\alpha:x\longrightarrow x'$, $\beta:y\longrightarrow y'$:
\[
\begin{tikzcd}[row sep=huge,column sep=huge]
 \varphi(x) \arrow[r,"{\varphi(\alpha)}"] \arrow[d,"{\gamma}"']
 & \varphi(x') \arrow[d,"{\gamma'}"]\\
 \psi(y) \arrow[r,"{\psi(\beta)}"'] & \psi(y')
\end{tikzcd}
\]

可用这一 $2$-范畴意义下的范畴纤维积来重新表述
(半)齐性余纤维群胚及光滑函子的概念.
设 $P:\underline A\longrightarrow\underline C_\Lambda$ 为余纤维范畴,
并考虑 $\underline C_\Lambda$ 中的 $R_1\times_RR_2$.
取正像函子
\[
 (\pi_i)_*:\underline A(R_1\times_RR_2)
 \longrightarrow\underline A(R_i),\qquad i=1,2.
\]
它们诱导函子
\[
 \underline A(R_1\times_RR_2)\longrightarrow
 \underline A(R_1)\times^2_{\underline A(R)}\underline A(R_2).
\]
在典范自然等价意义下, 此函子与所取的正像函子无关.
群胚 $\underline A$ 齐性当且仅当只要 $R_1\longrightarrow R$ 满射,
此函子就是等价; 它半齐性当且仅当
\[
 [\underline A(R_1\times_RR_2)]\longrightarrow
 [\underline A(R_1)]
 \times_{[\underline A(R)]}[\underline A(R_2)]
\]

%%%%%%%%%%  scan idx 65  |  folio 54  %%%%%%%%%%

在 $R_1\longrightarrow R$ 满射时为满射,
而在 $R=k$, $R_1=k[\varepsilon]$ 时为双射.

现设 $\varphi:\underline A\longrightarrow\underline B$
为 $\underline C_\Lambda$ 上余纤维群胚间的函子.
对 $\underline C_\Lambda$ 中每个映射 $R'\longrightarrow R$,
$\varphi$ 在典范自然等价意义下诱导函子
\[
 \underline A(R')\longrightarrow
 \underline A(R)\times^2_{\underline B(R)}\underline B(R').
\]
函子 $\varphi$ 光滑当且仅当只要 $R'\longrightarrow R$ 满射,
此函子就在对象的同构类上满射.

\underline{注 \SGARefTarget{sga7.sga7i.evi.n2-6.item-c}{2.6}(c).}
若 $\underline A$ 齐性, 则对任意有限维 $k$-向量空间 $V$ 和 $W$, 函子
\[
 \underline A(k[V\times W])\longrightarrow
 \underline A(k[V])\times\underline A(k[W])
\]
是范畴等价. 回顾约定 $\underline A(k)=\{e\}$. 特别地,
\[
 \operatorname{Aut}_{\underline A(k[V\times W])}(1_{V\times W})
 \xrightarrow{\ \sim\ }
 \operatorname{Aut}_{\underline A(k[V])}(1_V)
 \times
 \operatorname{Aut}_{\underline A(k[W])}(1_W),
\]
其中 $1_V$ 为 $e$ 在 $\underline C_\Lambda$ 中唯一映射
$k\longrightarrow k[V]$ 下的正像.
从而, 若 $\underline A$ 齐性, 则 $\operatorname{Aut}(1_\varepsilon)$ 及
$[\underline A(k[\varepsilon])]$ 均带有典范的 $k$-向量空间结构.

\underline{注 \SGARefTarget{sga7.sga7i.evi.n2-6.item-d}{2.6}(d).}
对每个 $R\in\widehat{\underline C}_\Lambda$,
$\underline C_\Lambda$ 上的余纤维离散群胚 $\underline h_R$ 齐性.

若 $\underline A$ 半齐性, 则 $[\underline A]$ 也半齐性.
但 $\underline A$ 的齐性一般不蕴含 $[\underline A]$ 的齐性.

\underline{命题 \SGARefTarget{sga7.sga7i.evi.n2-7}{2.7}.}
设 $\underline A$ 为 $\underline C_\Lambda$ 上的齐性余纤维群胚.
下列条件等价.

(i) $[\underline A]$ 齐性; 即对每个小扩张 $R'\longrightarrow R$,
\[
 [\underline A(R'\times_RS)]\longrightarrow
 [\underline A(R')]
 \times_{[\underline A(R)]}[\underline A(S)]
\]
为双射.

(ii) 对 $\underline A$ 中每个小延拓 $a'\longrightarrow a$, 映射
\[
 \operatorname{Aut}_{\underline A(R')}(a')
 \longrightarrow
 \operatorname{Aut}_{\underline A(R)}(a)
\]

%%%%%%%%%%  scan idx 66  |  folio 55  %%%%%%%%%%

满射, 其中 $R'=P(a')$, $R=P(a)$.

\underline{证明.} (i)$\Rightarrow$(ii).
设 $f:R'\longrightarrow R$ 为 $\underline C_\Lambda$ 中的小扩张,
$a'\longrightarrow a$ 为 $\underline A$ 中的 $f$-映射.
给定 $\tau\in\operatorname{Aut}_{\underline A(R)}(a)$,
考虑 $a'\times_a a'$ 及 $a'\times_a a'_\tau$,
其中 $a'_\tau\longrightarrow a$ 为复合
\[
 a'\longrightarrow a\xrightarrow{\ \tau\ }a.
\]
因
\[
 [\underline A(R'\times_RR')]
 \longrightarrow
 [\underline A(R')]
 \times_{[\underline A(R)]}[\underline A(R')]
\]
为双射, 存在同构
\[
 a'\times_a a'\xrightarrow{\ u\ }a'\times_a a'_\tau.
\]
再向第二个因子投影, 得到提升 $\tau$ 的同构 $\tau':a'\longrightarrow a'$.

(ii)$\Rightarrow$(i).
设 $a''_i$ $(i=1,2)$ 为 $\underline A(R'\times_RS)$ 的对象,
它们在映射
\[
 [\underline A(R'\times_RS)]\longrightarrow
 [\underline A(R')]
 \times_{[\underline A(R)]}[\underline A(S)]
\]
下有相同的像. 得到由实线箭头组成的图
\[
\begin{tikzcd}[column sep=small,row sep=small]
 && a''_2 \arrow[ddll] \arrow[ddrr] &&\\
 && a''_1 \arrow[u,dashed] \arrow[dl] \arrow[dr] &&\\
 a'_2 \arrow[ddrr] & a'_1 \arrow[l,"{\sim}"'] \arrow[dr]
 && b_1 \arrow[r,"{\tau}"] \arrow[dl] & b_2 \arrow[ddll]\\
 && a_1 \arrow[d,"{\bar\tau}"] &&\\
 && a_2 &&
\end{tikzcd}
\qquad\text{位于右图之上}\qquad
\begin{tikzcd}[column sep=small,row sep=small]
 & R'\times_RS \arrow[dl] \arrow[dr] &\\
 R' \arrow[dr] && S \arrow[dl]\\
 & R &
\end{tikzcd}
\]
其中 $\bar\tau:a_1\longrightarrow a_2$ 为 $\tau$ 诱导的同构,
故右下方方块交换. 由条件 (ii), 可以替换给定的同构
$a'_1\xrightarrow{\sim}a'_2$, 使左方方块也交换.
由注 \SGARefLink{sga7.sga7i.evi.n2-6.item-a}{2.6}(a),
\[
 a''_i=a'_i\times_{a_i}b_i\qquad(i=1,2)
\]
因而得到同构 $a''_1\xrightarrow{\sim}a''_2$. 故
\[
 [\underline A(R'\times_RS)]\longrightarrow
 [\underline A(R')]
 \times_{[\underline A(R)]}[\underline A(S)].
\]
为单射, 又因假定 $\underline A$ 齐性, 故为双射. \hfill$\square$

% END INCLUDED SGA7I_English_current_buildable_source_20260823/components/19_expose_VI_2_3_through_2_7.tex
% BEGIN INCLUDED SGA7I_English_current_buildable_source_20260823/components/20_expose_VI_2_8_through_2_9.tex
% R3 explicit revisions VI-A-02; see qa/FIDELITY-R3-CHANGES.json.
设 $\underline A$ 为 $\underline C_\Lambda$ 上的齐性余纤维群胚,
$f:R'\longrightarrow R$ 为 $\underline C_\Lambda$ 中的小扩张. 有典范映射

%%%%%%%%%%  scan idx 67  |  folio 56  %%%%%%%%%%

\[
 \widetilde f:R'\times_RR'\longrightarrow k[\operatorname{Ker}f],
 \qquad
 (r'_1,r'_2)\longmapsto r'_1(0)+(r'_1-r'_2),
\]
且
\[
 1\times\widetilde f:R'\times_RR'
 \longrightarrow R'\times_k k[\operatorname{Ker}f]
\]
为同构. 给定 $\underline A$ 中两个 $f$-映射
$a_i\longrightarrow a$ $(i=1,2)$, 纤维积 $a_1\times_a a_2$
为 $R'\times_RR'$ 上的对象. 因而可构造
\[
 \tau(a_1,a_2)=
 (\widetilde f)_*(a_1\times_a a_2)
 \in\underline A(k[\operatorname{Ker}f]).
\]
因 $1\times\widetilde f$ 为同构, 存在与到第一因子的投影相容的同构
\[
 a_1\times_a a_2\simeq a_1\times\tau(a_1,a_2)
\]
断言 \SGARefLink{sga7.sga7i.evi.n1-5}{1.5}(2) 现在具有下列更明确的形式.

\underline{引理 \SGARefTarget{sga7.sga7i.evi.n2-8}{2.8}.}
设 $\underline A$ 为 $\underline C_\Lambda$ 上的齐性余纤维群胚,
并设
\[
 a_i\xrightarrow{\ \alpha_i\ }a\qquad(i=1,2)
\]
为 $\underline A$ 中的 $f$-映射.

(1) 存在 $\underline A$ 中的 $\rho:a_1\longrightarrow a_2$
使 $\alpha_2\circ\rho=\alpha_1$, 当且仅当$\underline A$ 中存在箭头
$a_1\longrightarrow\tau(a_1,a_2)$.

(2) 存在满足 $P(\rho)=\operatorname{id}_{R'}$ 的这种 $\rho$, 当且仅当
\[
 \tau(a_1,a_2)=0
 \quad\text{在}\quad
 [\underline A(k[\operatorname{Ker}f])].
\]

\underline{证明.} 同时证明两个断言.
存在满足 $\alpha_2\circ\rho=\alpha_1$
(及 $P(\rho)=\operatorname{id}_{R'}$) 的 $\rho:a_1\longrightarrow a_2$,
依次等价于存在
\[
 h:a_1\longrightarrow a_1\times_a a_2,
 \qquad \operatorname{pr}_1\circ h=\operatorname{id}_{a_1},
\]
（对断言 (2), 还要求 $P(h)$ 为对角映射）; 等价于存在
\[
 g:a_1\longrightarrow a_1\times\tau(a_1,a_2),
 \qquad \operatorname{pr}_1\circ g=\operatorname{id}_{a_1},
\]
（对断言 (2), 还要求
\[
 P(g):R'\longrightarrow R'\times_k k[\operatorname{Ker}f],
 \qquad r'\longmapsto(r',r'(0));
\]
）
以及等价于 $\underline A$ 中存在箭头
$\gamma:a_1\longrightarrow\tau(a_1,a_2)$, （对断言 (2), 还要求
\[
 P(\gamma):R'\longrightarrow k[\operatorname{Ker}f],
 \qquad r'\longmapsto r'(0).
\]
）
这样的箭头经过 $R'\longrightarrow k\longrightarrow k[\operatorname{Ker}f]$
分解当且仅当 $\tau(a_1,a_2)=0$. \hfill$\square$

\underline{命题 \SGARefTarget{sga7.sga7i.evi.n2-9}{2.9}.} (a) 设
\[
 \underline A\xrightarrow{\ \varphi\ }\underline B
 \xrightarrow{\ \psi\ }\underline C
\]
为 $\underline C_\Lambda$ 上余纤维群胚间的函子. 假定

(i) $\underline B$ 半齐性而 $\underline C$ 齐性;

(ii) 映射
\[
 \psi(k[\varepsilon]):[\underline B(k[\varepsilon])]
 \longrightarrow[\underline C(k[\varepsilon])]
\]
为单射.

则 $\psi\circ\varphi$ 光滑当且仅当 $\varphi$ 和 $\psi$ 均

%%%%%%%%%%  scan idx 68  |  folio 57  %%%%%%%%%%

光滑.

(b) 设 $\varphi:\underline A\longrightarrow\underline B$
为 $\underline C_\Lambda$ 上余纤维群胚间的光滑函子.
若 $\underline A$ 半齐性而 $\underline B$ 齐性, 则
\[
 \widehat\varphi:\widehat{\underline A}
 \longrightarrow\widehat{\underline B}
\]
光滑. 明确地说, 对 $\widehat{\underline B}$ 中每个满射
$b'\xrightarrow{\beta}\widehat\varphi(a)$, 存在
$\widehat{\underline A}$ 中的映射 $\alpha:a'\longrightarrow a$
及 $\widehat{\underline C}_\Lambda$-同构
\[
 \rho:\widehat\varphi(a')\xrightarrow{\ \sim\ }b'
\]
使 $\beta\circ\rho=\varphi(\alpha)$.

\underline{证明.} (a) 假定 $\psi\circ\varphi$ 光滑.
若 $\varphi$ 光滑, 则由命题 \SGARefLink{sga7.sga7i.evi.n2-3}{2.3}(b),
$\psi$ 光滑, 故只须证明 $\varphi$ 光滑.
设 $b'\xrightarrow{\beta}\varphi(a)$ 为 $\underline B$ 中的小延拓.
因 $\psi\circ\varphi$ 光滑, 存在 $\underline A$ 中的箭头
$a'\xrightarrow{\alpha}a$ 及 $\underline C_\Lambda$-同构
\[
 \rho:\psi(b')\longrightarrow\psi\varphi(a')
\]
使图
\[
\begin{tikzcd}[column sep=huge,row sep=huge]
 \psi(b') \arrow[rr,"{\rho}"] \arrow[dr,"{\psi(\beta)}"']
 && \psi\varphi(a') \arrow[dl,"{\psi\varphi(\alpha)}"]\\
 & \psi\varphi(a) &
\end{tikzcd}
\]
交换. 因 $\underline B$ 半齐性, 可取交换图
\[
\begin{tikzcd}[column sep=small,row sep=small]
 & b'' \arrow[dl] \arrow[dr] &
 &&& R'\times_RR' \arrow[dl] \arrow[dr] &\\
 b' \arrow[dr,"{\beta}"'] &&
 \varphi(a') \arrow[dl,"{\varphi(\alpha)}"]
 & \text{位于右图之上} &
 R' \arrow[dr] && R' \arrow[dl]\\
 & \varphi(a) &&&& R &
\end{tikzcd}
\]
其中 $R=P(a)$, $R'=P(b')=P(a')$.
由 $\underline C$ 的齐性及注 \SGARefLink{sga7.sga7i.evi.n2-6.item-a}{2.6}(a),
\[
 \psi(b'')=
 \psi(b')\times_{\psi\varphi(a)}\psi\varphi(a')
\]
引理 \SGARefLink{sga7.sga7i.evi.n2-8}{2.8}(2) 于是给出
\[
 0=\tau(\psi\varphi(a'),\psi(b'))
   =(\widetilde f)_*\psi(b'').
\]
由假设 (ii), 得 $(\widetilde f)_*(b'')=0$.
断言 \SGARefLink{sga7.sga7i.evi.n1-5}{1.5}(2) 现在给出映射
$\gamma:b'\longrightarrow\varphi(a')$ 满足
$\beta=\varphi(\alpha)\circ\gamma$.
因此 $\varphi$ 光滑, 从而 $\psi$ 也光滑.
逆向由命题 \SGARefLink{sga7.sga7i.evi.n2-3}{2.3}(a) 得到.

(b) 给定 $\widehat{\underline B}$ 中的满射
$\beta:b'\longrightarrow\widehat\varphi(a)$.
它由 $\underline B$ 中的交换图序列表示:

%%%%%%%%%%  scan idx 69  |  folio 58  %%%%%%%%%%

\[
\begin{tikzcd}[column sep=large]
 {} \arrow[r,dashed] & b'_3 \arrow[r,"{\beta'_2}"] \arrow[d,"{\beta_3}"]
 & b'_2 \arrow[r,"{\beta'_1}"] \arrow[d,"{\beta_2}"]
 & b'_1 \arrow[d,"{\beta_1}"]\\
 {} \arrow[r,dashed] & \varphi(a_3) \arrow[r,"{\varphi(\alpha_2)}"']
 & \varphi(a_2) \arrow[r,"{\varphi(\alpha_1)}"']
 & \varphi(a_1)
\end{tikzcd}
\]
其中每个 $\alpha_v$、$\beta_v$、$\beta'_v$ 都满射,
且 $\beta=\varprojlim_v\beta_v$.
只须归纳证明下述每个由实线箭头组成的图,
若 $\alpha_n$、$\beta_n$、$\beta_{n+1}$、$\beta'_n$ 满射,
均能添加点线箭头补全:
\[
\begin{tikzcd}[column sep=large,row sep=large]
 \varphi(a'_{n+1})
   \arrow[rrrr,dashed,"{\varphi(\,)}"]
   \arrow[dr,dashed,leftrightarrow]
   \arrow[ddr,dashed,"{\varphi(\,)}"']
 &&&& \varphi(a'_n)
   \arrow[dl,leftrightarrow]
   \arrow[ddl,"{\varphi(\alpha'_n)}"]\\
 & b'_{n+1} \arrow[rr,"{\beta'_n}"] \arrow[d,"{\beta_{n+1}}"']
 && b'_n \arrow[d,"{\beta_n}"'] &\\
 & \varphi(a_{n+1}) \arrow[rr,"{\varphi(\alpha_n)}"']
 && \varphi(a_n) &
\end{tikzcd}
\]
由假设, $\underline A$ 半齐性且 $\alpha_n:a'_{n}\longrightarrow a_n$ 满射.
故存在交换图
\[
\begin{tikzcd}[column sep=0.8em,row sep=small]
 & a' \arrow[dl,"{\gamma}"'] \arrow[dr,"{\gamma'}"] &
 &&& R_{n+1}\times_{R_n}R'_n \arrow[dl] \arrow[dr] &\\
 a_{n+1} \arrow[dr,"{\alpha_n}"'] &&
 a'_n \arrow[dl,"{\alpha'_n}"]
 & \text{在 $\underline A$ 中位于右图之上} &
 R_{n+1} \arrow[dr] && R'_n \arrow[dl]
 & \text{在 $\underline C_\Lambda$ 中}\\
 & a_n &&&& R_n &&
\end{tikzcd}
\]
$\underline B$ 的齐性蕴含
\[
\begin{tikzcd}[column sep=huge,row sep=huge]
 & \varphi(a') \arrow[dl,"{\varphi(\gamma)}"']
                  \arrow[dr,"{\varphi(\gamma')}"] &\\
 \varphi(a_{n+1}) \arrow[dr,"{\varphi(\alpha_n)}"']
 && \varphi(a'_n) \arrow[dl,"{\varphi(\alpha'_n)}"]\\
 & \varphi(a_n) &
\end{tikzcd}
\]

%%%%%%%%%%  scan idx 70  |  folio 59  %%%%%%%%%%

为 $\underline B$ 中的纤维积图.
因此存在满足所需交换性的唯一映射
$\delta:b'_{n+1}\longrightarrow\varphi(a')$.
因 $\beta_{n+1}$、$\beta'_n$ 满射, $\delta$ 也满射.
故 $\varphi$ 的光滑性给出 $\underline A$ 中的映射
$a'_{n+1}\xrightarrow{\alpha'}a'$ 及 $\underline C_\Lambda$-同构
\[
 \rho:\varphi(a'_{n+1})\xrightarrow{\ \sim\ }b'_{n+1}
\]
使 $\delta\circ\rho=\varphi(\alpha')$. 得到交换图
\[
\begin{tikzcd}[row sep=huge,column sep=huge]
 \varphi(a'_{n+1})
   \arrow[rrrr,"{\varphi(\gamma'\circ\alpha')}"]
   \arrow[dr,"{\rho}"]
   \arrow[ddr,"{\varphi(\gamma\circ\alpha')}"']
 &&&& \varphi(a'_n)
   \arrow[dl,"{\sim}"']
   \arrow[ddl,"{\varphi(\alpha'_n)}"]\\
 & b'_{n+1} \arrow[rr] \arrow[d]
 && b'_n \arrow[d] &\\
 & \varphi(a_{n+1}) \arrow[rr]
 && \varphi(a_n) &
\end{tikzcd}
\]
这就完成归纳并证明了 (b). \hfill$\square$

% END INCLUDED SGA7I_English_current_buildable_source_20260823/components/20_expose_VI_2_8_through_2_9.tex

% BEGIN INCLUDED SGA7I_English_current_buildable_source_20260823/components/21_expose_VI_corollary_2_10.tex
\underline{推论 \SGARefTarget{sga7.sga7i.evi.n2-10}{2.10}.}
(a) 设 $\varphi:\underline A\longrightarrow\underline B$
为 $\underline C_\Lambda$ 上余纤维群胚间的光滑函子.
设 $\xi$ 为 $\underline A$ 的遍有形式对象,
$\eta$ 为 $\underline B$ 的极小遍有形式对象.
若 $\underline B$ 齐性, 则对 $\widehat{\underline B}$ 中每个映射
$\beta:\eta\longrightarrow\varphi(\xi)$,
环 $P(\xi)$ 经由
\[
 P(\beta):P(\eta)\longrightarrow P(\varphi(\xi))=P(\xi).
\]
成为 $P(\eta)$ 上的形式幂级数环.

(b) 设 $\underline A$ 为 $\underline C_\Lambda$ 上的齐性余纤维群胚,
$\xi$ 和 $\eta$ 为 $\underline A$ 的遍有形式对象.
若 $\eta$ 极小, 则对 $\widehat{\underline A}$ 中每个映射
$\alpha:\eta\longrightarrow\xi$, 环 $P(\xi)$ 经由
\[
 P(\alpha):P(\eta)\longrightarrow P(\xi).
\]
成为 $P(\eta)$ 上的形式幂级数环.

\underline{证明.} (a) 映射 $\beta:\eta\longrightarrow\varphi(\xi)$ 诱导交换图
\[
\begin{tikzcd}[row sep=huge,column sep=huge]
 \underline h_{P(\xi)} \arrow[r,"{\widetilde\beta}"]
                           \arrow[d,"{\widetilde\xi}"']
 & \underline h_{P(\eta)} \arrow[d,"{\widetilde\eta}"]\\
 \underline A \arrow[r,"{\varphi}"'] & \underline B.
\end{tikzcd}
\]
因 $\widetilde\xi$ 和 $\varphi$ 光滑, 复合
$\widetilde\eta\circ\widetilde\beta=\varphi\circ\widetilde\xi$ 光滑.
另一方面,

%%%%%%%%%%  scan idx 71  |  folio 60  %%%%%%%%%%

\[
 \widetilde\eta(k[\varepsilon]):
 \underline h_{P(\eta)}(k[\varepsilon])
 \longrightarrow[\underline B(k[\varepsilon])]
\]
为同构. 因此命题 \SGARefLink{sga7.sga7i.evi.n2-9}{2.9}(a)
表明 $\widetilde\beta$ 光滑, 恰好就是说 $P(\xi)$ 是 $P(\eta)$ 上的形式幂级数环.

(b) 对 $\underline A=\underline B$ 应用 (a). \hfill$\square$

% END INCLUDED SGA7I_English_current_buildable_source_20260823/components/21_expose_VI_corollary_2_10.tex
% BEGIN INCLUDED SGA7I_English_current_buildable_source_20260823/components/22_expose_VI_representable_cogroupoids_intro.tex
最后考虑余纤维群胚的“可表性”. 设 $\underline C$ 为固定范畴.
$\underline C$ 的每个对象 $R$ 定义函子
\[
 h_R:\underline C\longrightarrow(\mathrm{Sets}),
 \qquad h_R(S)=\operatorname{Hom}_{\underline C}(R,S),
\]
从而得到全忠实函子
\[
 h:\underline C^{\mathrm{op}}
 \longrightarrow
 \underline{\operatorname{Hom}}(\underline C,(\mathrm{Sets})).
\]
\emph{$\underline C$ 中的余范畴(cocategory)}是图
\[
\begin{tikzcd}[column sep=huge]
 h_{R_1}\times_{h_{R_0}}h_{R_1} \arrow[r]
 & h_{R_1}
   \arrow[r,shift left=3,"{s_1}"]
   \arrow[r,shift right=3,"{s_2}"']
 & h_{R_0} \arrow[l,"{\epsilon}" description]
\end{tikzcd}
\qquad
\text{满足 }s_i\circ\epsilon=\operatorname{id}\quad(i=1,2),
\]
其中首个箭头记为 $c$, 且
\[
 h_{R_1}\times_{h_{R_0}}h_{R_1}
 = (h_{R_1})_{s_1}\times_{h_{R_0}}{}_{s_2}(h_{R_1}),
\]
满足

(i) $c$ 有结合性, 并与 $s_i$ ($i=1,2$) 相容;

(ii) 两个复合
\begin{align*}
 h_{R_1}&\xrightarrow{\ \sim\ }
 h_{R_0}\times_{h_{R_0}}h_{R_1}
 \xrightarrow{\ \epsilon\times1\ }
 h_{R_1}\times_{h_{R_0}}h_{R_1}
 \xrightarrow{\ c\ }h_{R_1},\\[1ex]
 h_{R_1}&\xrightarrow{\ \sim\ }
 h_{R_1}\times_{h_{R_0}}h_{R_0}
 \xrightarrow{\ 1\times\epsilon\ }
 h_{R_1}\times_{h_{R_0}}h_{R_1}
 \xrightarrow{\ c\ }h_{R_1}
\end{align*}
均为恒等映射.

若 $(R_0,R_1,s_1,s_2,\epsilon,c)$ 为 $\underline C$ 中的余范畴,
则对每个 $S\in\underline C$, 定义范畴 $h_{(R_0,R_1,\ldots)}(S)$ 如下:
\[
 \operatorname{ob}h_{(R_0,R_1,\ldots)}(S)
   =\operatorname{Hom}_{\underline C}(R_0,S),
 \qquad
 \operatorname{Fl}h_{(R_0,R_1,\ldots)}(S)
   =\operatorname{Hom}_{\underline C}(R_1,S),
\]
其复合由
\[
 c:h_{R_1}\times_{h_{R_0}}h_{R_1}\longrightarrow h_{R_1}.
\]
给出. 所得范畴函子性地依赖于 $S$,
并定义 $\underline C$ 上的余纤维范畴 $\underline h_{(R_0,R_1,\ldots)}$.
例如 $\underline C$ 的每个对象 $R$ 定义离散余范畴
\[
 \underline h_R
 =\underline h_{(R,R,s_1,s_2,\epsilon,c)}
\]
只须取 $s_1,s_2,\epsilon$ 及 $c$ 均为恒等映射. 余范畴的同态
\[
 (R_0,R_1,\ldots)\longrightarrow(R'_0,R'_1,\ldots)
\]

%%%%%%%%%%  scan idx 72  |  folio 61  %%%%%%%%%%

为一对可表函子间的映射 $h_{R_0}\longrightarrow h_{R'_0}$ 及
$h_{R_1}\longrightarrow h_{R'_1}$,
使其诱导的映射
\[
 h_{(R_0,R_1,\ldots)}(S)
 \longrightarrow h_{(R'_0,R'_1,\ldots)}(S)
\]
对每个 $S\in\underline C$ 都是函子.
特别地, 每个余范畴 $(R_0,R_1,\ldots)$ 都有同态
\[
 \epsilon:\underline h_{R_0}
 \longrightarrow\underline h_{(R_0,R_1,\ldots)}.
\]
\emph{$\underline C$ 中的余群胚(cogroupoid)}为余范畴 $(R_0,R_1,\ldots)$,
使对每个 $S\in\underline C$, $h_{(R_0,R_1,\ldots)}(S)$ 都是群胚;
等价地, $\underline h_{(R_0,R_1,\ldots)}$ 是 $\underline C$
上以群胚为纤维的余纤维范畴.

实际中, 余范畴与余群胚如下出现.
设 $P:\underline F\longrightarrow\underline C$ 为余纤维范畴,
固定 $\underline F$ 的对象 $\xi$.
置 $R_0=P(\xi)$, 并以 $(R_0,\underline C)$ 表示
$\underline C$ 中源为 $R_0$ 的箭头范畴.
设 $P:\underline F'\longrightarrow(R_0,\underline C)$ 为
$\underline F$ 沿典范函子 $(R_0,\underline C)\longrightarrow\underline C$ 的拉回.
取 $\underline F'$ 在 $(R_0,\underline C)$ 上的余笛卡尔截面 $\xi^{\#}$,
使 $\xi^{\#}(\operatorname{id}_{R_0})=\xi$.
即对 $\underline C$ 中每个箭头 $R_0\xrightarrow fS$,
取一个余笛卡尔的 $f$-映射 $\xi\longrightarrow\xi^{\#}(f)$.

定义
\[
 \operatorname{Hom}_{\xi^{\#}}:
 (R_0\amalg R_0,\underline C)\longrightarrow(\mathrm{Sets})
\]
为
\[
 X=(f_1,f_2)\longmapsto
 \operatorname{Hom}_{\underline F(S)}
   (\xi^{\#}(f_1),\xi^{\#}(f_2)),
\]
其中 $S$ 为 $f_1$ 和 $f_2$ 的共同靶.
在典范自然等价意义下, 此函子与 $\xi^{\#}$ 的选择无关;
此后将 $\xi^{\#}$ 简记为 $\xi$.

假定此函子由 $(R_0\amalg R_0,\underline C)$ 中的对象
\[
 X_0=
 \left(R_0\mathrel{\substack{\xrightarrow{\ s_1\ }\\[-.6ex]
                              \xrightarrow[\ s_2\ ]{}}}R_1\right)
\]
表示. 即在 $\underline F$ 中存在 $R_0$-映射
$\varphi:\xi(s_1)\longrightarrow\xi(s_2)$, 使
\[
 \widetilde\varphi:
 \operatorname{Hom}_{(R_0\amalg R_0,\underline C)}(X_0,X)
 \longrightarrow\operatorname{Hom}_\xi(X)
\]
对 $(R_0\amalg R_0,\underline C)$ 中每个 $X$ 都是双射.
则 $X_0$ 自然定义 $\underline C$ 中的余范畴 $(R_0,R_1,\ldots)$.
事实上, $\underline C$ 中的三个箭头

%%%%%%%%%%  scan idx 73  |  folio 62  %%%%%%%%%%

\[
\begin{tikzcd}[column sep=huge]
 R_0 \arrow[r,shift left=3,"u"]
     \arrow[r,"v"]
     \arrow[r,shift right=3,"w"']
 & S
\end{tikzcd}
\]
定义复合
\[
 \operatorname{Hom}_{\underline F(S)}(\xi(v),\xi(w))
 \times
 \operatorname{Hom}_{\underline F(S)}(\xi(u),\xi(v))
 \longrightarrow
 \operatorname{Hom}_{\underline F(S)}(\xi(u),\xi(w)),
\]
从而给出关于 $S$ 函子性的合成法则
\[
 h_{R_1}(S)\times_{h_{R_0}(S)}h_{R_1}(S)
 \longrightarrow h_{R_1}(S),
\]
因此 $\operatorname{Hom}_\xi$ 定义 $\underline C$ 中的余范畴
$(R_0,R_1,\ldots)$ 及 $\underline C$ 上的余纤维范畴
$\underline h_{(R_0,R_1,\ldots)}$. 明确地说,
\[
 \operatorname{ob}\underline h_{(R_0,R_1,\ldots)}
 =\coprod_{S\in\operatorname{ob}(\underline C)}
   \operatorname{Hom}_{\underline C}(R_0,S).
\]
箭头
\[
 (R_0\xrightarrow fS)\longrightarrow(R_0\xrightarrow{f'}S')
\]
为 $\underline C$ 中一对箭头 $S\xrightarrow hS'$ 和 $R_1\xrightarrow uS'$,
使
\[
 u\circ s_1=h\circ f,
 \qquad u\circ s_2=f'.
\]
结构函子 $P:\underline h_{(R_0,R_1,\ldots)}\longrightarrow\underline C$
将 $(R_0\xrightarrow fS)$ 送到 $S$.

二元组 $(\xi,\varphi)$ 现定义函子
\[
 (\xi,\varphi):\underline h_{(R_0,R_1,\ldots)}
 \longrightarrow\underline F
 \quad\text{在 }\underline C
\]
其在对象上为
\[
 (R_0\xrightarrow fS)\longmapsto\xi(f)
\]
在箭头上为
\[
 \left[(R_0\xrightarrow fS)
 \xrightarrow{(h,u)}(R_0\xrightarrow{f'}S')\right]
 \longmapsto
 \left[\xi(f)\xrightarrow{h_*}\xi(hf)
 \xrightarrow{u(\varphi)}\xi(f')\right]
\]
对每个 $S\in\underline C$, 由 $\operatorname{Hom}_\xi$ 的可表性, 函子
\[
 (\xi,\varphi)(S):
 \underline h_{(R_0,R_1,\ldots)}(S)\longrightarrow\underline F(S)
\]
全忠实. 故
\[
 (\xi,\varphi):\underline h_{(R_0,R_1,\ldots)}
 \longrightarrow\underline F
\]
为 $\underline C$-等价当且仅当 $\xi$ 是 $\underline F$
在 $\underline C$ 上的拟始对象, 即 $\underline F$ 的每个对象 $a$
都容许余笛卡尔箭头 $\xi\longrightarrow a$.

也可以改为考虑函子
\[
 \operatorname{Isom}_\xi:
 (R_0\amalg R_0,\underline C)\longrightarrow(\mathrm{Sets}),
\]

%%%%%%%%%%  scan idx 74  |  folio 63  %%%%%%%%%%

由
\[
 X=(f_1,f_2)\longmapsto
 \operatorname{Isom}_{\underline F(S)}(\xi(f_1),\xi(f_2)),
\]
给出, 其中 $S$ 为 $f_1$ 和 $f_2$ 的共同靶.
若此函子由 $(R_0\amalg R_0,\underline C)$ 中的对象
$(R_0\rightrightarrows R_1)$ 表示,
则它定义 $\underline C$ 中的余群胚 $(R_0,R_1,\ldots)$ 及函子
\[
 (\xi,\varphi):\underline h_{(R_0,R_1,\ldots)}
 \longrightarrow\underline F.
\]
当 $\underline F$ 为余纤维群胚时, 此函子为 $\underline C$-等价当且仅当
$\xi$ 是 $\underline F$ 在 $\underline C$ 上的拟始对象.

% END INCLUDED SGA7I_English_current_buildable_source_20260823/components/22_expose_VI_representable_cogroupoids_intro.tex

% BEGIN INCLUDED SGA7I_English_current_buildable_source_20260823/components/23_expose_VI_2_11_and_2_12_ab.tex
\underline{定义与命题 \SGARefTarget{sga7.sga7i.evi.n2-11}{2.11}.}
若下列等价条件之一(因而二者)成立, 则称
$\widehat{\underline C}_\Lambda$ 中的余群胚 $(R_0,R_1,\ldots)$ 为\underline{光滑的}:

(i) $R_1$ 经由 $s_1$ (以及经由 $s_2$) 是 $R_0$ 上的形式幂级数环;

(ii) 函子
\[
 \epsilon:\underline h_{R_0}
 \longrightarrow\underline h_{(R_0,R_1,\ldots)}
\]
在 $\underline C_\Lambda$ 上光滑.

$\widehat{\underline C}_\Lambda$ 中的\underline{极小光滑余群胚},
也称\underline{正规化光滑余群胚}, 是光滑余群胚 $(R_0,R_1,\ldots)$,
满足 $\underline h_{(R_0,R_1,\ldots)}(k[\varepsilon])$ 完全不连通;
等价地,
\[
 s_1(k[\varepsilon])=s_2(k[\varepsilon]).
\]

\underline{证明.}
$\epsilon:\underline h_{R_0}\to\underline h_{(R_0,R_1,\ldots)}$
光滑就是说, 对 $\underline h_{(R_0,R_1,\ldots)}$ 中每个满射
\[
 (S',f')\xrightarrow{(h,u)}(S,f)
\]
存在 $f$ 到 $S'$ 的提升 $f''$ 及同构
\[
 (S',f'')\xrightarrow{(1,u')}(S',f').
\]
等价地, 对 $\underline C_\Lambda$ 中每个满射 $S'\to S$,
给定 $u\in h_{R_1}(S)$ 及 $s_1(u)$ 到 $S'$ 的提升 $f'$,
存在 $u$ 到 $S'$ 的提升 $u'$ 满足 $s_1(u')=f'$.
这正是
\[
 s_1:\underline h_{R_1}\longrightarrow\underline h_{R_0}
\]
在 $\underline C_\Lambda$ 上的光滑性, 又等价于
$R_1$ 经由 $s_1$ 是 $R_0$ 上的形式幂级数环. \hfill$\square$

\underline{命题 \SGARefTarget{sga7.sga7i.evi.n2-12}{2.12}.}
设 $(R_0,R_1,\ldots)$ 为 $\widehat{\underline C}_\Lambda$ 中的余群胚.

(a) 若 $(R_0,R_1,\ldots)$ 光滑, 则 $\underline h_{(R_0,R_1,\ldots)}$ 齐性.

(b) 若 $s_1(k[\varepsilon])=s_2(k[\varepsilon])$, 则下列条件等价:

(i) $(R_0,R_1,\ldots)$ 光滑;

(ii) $\underline h_{(R_0,R_1,\ldots)}$ 齐性;

(iii) $\underline h_{(R_0,R_1,\ldots)}$ 半齐性.

%%%%%%%%%%  scan idx 75  |  folio 64  %%%%%%%%%%

(c) 若 $\underline h_{(R_0,R_1,\ldots)}$ 半齐性,
则它与 $\underline h_{(S_0,S_1,\ldots)}$ 等价,
其中 $S_*$ 为某个极小光滑余群胚.

\underline{证明.} (a) 考虑 $\underline h_{(R_0,R_1,\ldots)}$ 中的图
\[
\begin{tikzcd}[column sep=huge,row sep=large]
 (S_1,f_1) \arrow[dr,"{(h_1,u_1)}"']
 && (S_2,f_2) \arrow[dl,"{(h_2,u_2)}"]\\
 & (S,f) &
\end{tikzcd}
\]
其中 $h_2:S_2\longrightarrow S$ 在 $\underline C_\Lambda$ 中满射.
因 $s_1:\underline h_{R_1}\longrightarrow\underline h_{R_0}$ 光滑,
$u_1^{-1}u_2$ 可提升为元素 $v\in h_{R_1}(S_2)$, 满足 $s_1(v)=f_2$.
有
\[
 h_2\circ s_2(v)=h_1\circ f_1,
\]
故 $f_1\times s_2(v)$ 是 $h_{R_0}(S_1\times_SS_2)$ 的元素, 并给出交换图
\[
\begin{tikzcd}[column sep=huge,row sep=large]
 & (S_1\times_SS_2,,f_1\times s_2(v))
   \arrow[dl,"{\pi_1}"']
   \arrow[dr,"{(\pi_2,v^{-1})}"] &\\
 (S_1,f_1) \arrow[dr,"{(h_1,u_1)}"']
 && (S_2,f_2) \arrow[dl,"{(h_2,u_2)}"]\\
 & (S,f) &
\end{tikzcd}
\]
我们断言, 其顶端对象为 $\underline h_{(R_0,R_1,\ldots)}$ 中的纤维积
\[
 (S_1,f_1)\times_{(S,f)}(S_2,f_2)
\]
事实上, 对每个对象 $(T,g)$, 到顶端对象的映射等价于满足
\begin{align*}
 h_1w_1&=h_2w_2,&
 s_1(w_1)&=h'_1g,&s_2(w_1)&=f_1,\\
 s_1(w_2)&=h'_2g,&
 s_2(w_2)&=s_2(v).
\end{align*}
的元组 $(h'_1\times h'_2,w_1,w_2)$.
置 $w'_2=v^{-1}w_2$, 得到它们与满足
\begin{align*}
 u_2h_2w'_2&=u_1h_1w_1,&
 s_1(w_1)&=h'_1g,&s_2(w_1)&=f_1,\\
 s_1(w'_2)&=h'_2g,&
 s_2(w'_2)&=f_2.
\end{align*}
的元组 $(h'_1\times h'_2,w_1,w'_2)$ 之间的双射.
后者恰为从 $(T,g)$ 到 $(S_1,f_1)$ 及 $(S_2,f_2)$ 的一对箭头,
其到 $(S,f)$ 的复合相同. 因此
\begin{multline*}
 \operatorname{Hom}((T,g),(S_1\times_SS_2,
                 f_1\times s_2(v)))\\
 \simeq
 \operatorname{Hom}((T,g),(S_1,f_1))
 \times_{\operatorname{Hom}((T,g),(S,f))}
 \operatorname{Hom}((T,g),(S_2,f_2)),
\end{multline*}
这证明了 (a).

%%%%%%%%%%  scan idx 76  |  folio 65  %%%%%%%%%%

(b) 只须证明 (iii)$\Rightarrow$(i).
若 $\underline h_{(R_0,R_1,\ldots)}$ 半齐性, 则
\[
 \dim[\underline h_{(R_0,R_1,\ldots)}(k[\varepsilon])]
 =\dim h_{R_0}(k[\varepsilon])<\infty,
\]
故它有极小遍有形式对象. 另一方面,
$(R_0,1_{R_0})$ 是 $\widehat{\underline h}_{(R_0,R_1,\ldots)}$
的极小拟始对象, 因为
\[
 h_{R_0}(k[\varepsilon])
 \longrightarrow
 [\underline h_{(R_0,R_1,\ldots)}(k[\varepsilon])]
\]
为双射. 因此由命题 \SGARefLink{sga7.sga7i.evi.n1-10}{1.10}(ii),
$(R_0,1_{R_0})$ 是 $\underline h_{(R_0,R_1,\ldots)}$ 的极小遍有形式对象.
等价地,
\[
 \epsilon:\underline h_{R_0}
 \longrightarrow\underline h_{(R_0,R_1,\ldots)}
\]
极小光滑. 这证明了 (i).

% END INCLUDED SGA7I_English_current_buildable_source_20260823/components/23_expose_VI_2_11_and_2_12_ab.tex
% BEGIN INCLUDED SGA7I_English_current_buildable_source_20260823/components/24_expose_VI_2_12_c_through_2_14.tex
(c) 设 $(S_0,\xi_0)$ 为 $\underline h_{(R_0,R_1,\ldots)}$
的极小遍有形式对象(定理 \SGARefLink{sga7.sga7i.evi.n1-11}{1.11}).
对任意对象
\[
 \xi:R_0\longrightarrow T
 \qquad\text{属于 }\underline h_{(R_0,R_1,\ldots)}(T),
\]
考虑函子
\[
 \operatorname{Isom}(\xi,\xi):
 \underline C_\Lambda/(T\widehat\otimes_\Lambda T)
 \longrightarrow(\mathrm{Sets}).
\]
它将 $R$ 对应到由
\[
 R_0\xrightarrow{\ \xi\ }T
 \mathrel{\substack{\xrightarrow{}\\[-.6ex]\xrightarrow[]{}}}
 T\widehat\otimes_\Lambda T\longrightarrow R.
\]
得到的 $\xi$ 在 $R$ 上的两个像之间的同构集合.
此函子由
\[
 R_1\otimes_{R_0\widehat\otimes_\Lambda R_0}
       (T\widehat\otimes_\Lambda T).
\]
表示. 取逆极限, 可见 $\operatorname{Isom}(\xi_0,\xi_0)$
由一个对象 $S_1$ 前表示, 且
\[
 \underline h_{(R_0,R_1,\ldots)}
 \simeq\underline h_{(S_0,S_1,\ldots)}.
\]
由 (b), 余群胚 $S_*$ 极小光滑. \hfill$\square$

\underline{定义 \SGARefTarget{sga7.sga7i.evi.n2-13}{2.13}.}
称 $\underline C_\Lambda$ 上的余纤维群胚 $\underline A$
为\underline{(极小)前可表的}, 若它在 $\underline C_\Lambda$-等价意义下
等同于 $\underline h_{(R_0,R_1,\ldots)}$,
其中 $(R_0,R_1,\ldots)$ 为 $\widehat{\underline C}_\Lambda$
中的某个(正规化)余群胚.

\underline{命题 \SGARefTarget{sga7.sga7i.evi.n2-14}{2.14}.}
若 $(R_0,R_1,\ldots)$ 为 $\widehat{\underline C}_\Lambda$
中的正规化余群胚, 则每个同态
\[
 (\alpha_0,\alpha_1):(R_0,R_1,\ldots)
 \longrightarrow(R_0,R_1,\ldots)
\]
只要诱导 $\underline C_\Lambda$-等价
\[
 \underline h_{(\alpha_0,\alpha_1)}:
 \underline h_{(R_0,R_1,\ldots)}
 \longrightarrow\underline h_{(R_0,R_1,\ldots)}
\]
就是同构. 特别地,

%%%%%%%%%%  scan idx 77  |  folio 66  %%%%%%%%%%

$\underline C_\Lambda$ 上余纤维群胚的正规化前表示若存在, 则在同构意义下唯一.

\underline{证明.} 考虑交换图
\[
\begin{tikzcd}[column sep=huge,row sep=huge]
 \underline h_{(R_0,R_1,\ldots)}
   \arrow[r,"{\underline h_{(\alpha_0,\alpha_1)}}"]
 & \underline h_{(R_0,R_1,\ldots)}\\
 \underline h_{R_0} \arrow[u,"\epsilon"']
   \arrow[r,"{\underline h_{\alpha_0}}"']
 & \underline h_{R_0} \arrow[u,"\epsilon"]
\end{tikzcd}
\]
其中 $\underline h_{(\alpha_0,\alpha_1)}$ 是 $\underline C_\Lambda$-等价.
因 $\underline h_{(R_0,R_1,\ldots)}(k[\varepsilon])$ 完全不连通,
\[
 [\underline h_{(\alpha_0,\alpha_1)}(k[\varepsilon])]
 \quad\text{及}\quad
 [\epsilon(k[\varepsilon])]
\]
均为双射, 因而 $\underline h_{\alpha_0}(k[\varepsilon])$ 也为双射.
故 $\alpha_0:R_0\longrightarrow R_0$ 是满射自态射, 因而是自同构.
随之 $\alpha_1$ 也为 $R_1$ 的自同构. \hfill$\square$

% END INCLUDED SGA7I_English_current_buildable_source_20260823/components/24_expose_VI_2_12_c_through_2_14.tex

% BEGIN INCLUDED SGA7I_English_current_buildable_source_20260823/components/25_expose_VI_2_15_through_2_17.tex
\underline{引理 \SGARefTarget{sga7.sga7i.evi.n2-15}{2.15}.}
设 $A$ 和 $B$ 为 $\underline C_\Lambda$ 上的离散群胚,
$\varphi:A\longrightarrow B$ 为极小光滑函子.
若 $A$ 和 $B$ 均齐性, 则 $\varphi$ 为 $\underline C_\Lambda$-等价.

\underline{证明.} 须证明对每个 $S\in\operatorname{ob}(\underline C_\Lambda)$,
$\varphi(S):A(S)\longrightarrow B(S)$ 均为双射.
对每个 $S$ 它为满射, 对 $S=k[\varepsilon]$ 它为双射.
对 $S$ 的长度用归纳法. 设 $S'\xrightarrow fS$ 为 $\underline C_\Lambda$ 中的小扩张.
因 $A$ 和 $B$ 为齐性离散群胚, 典范同构
\[
 1\times\widetilde f:S'\times_SS'
 \xrightarrow{\ \sim\ }S'\times_k k[\operatorname{Ker}f]
\]
诱导交换图

%%%%%%%%%%  scan idx 78  |  folio 67  %%%%%%%%%%

\[
\begin{tikzcd}[column sep=huge,row sep=huge]
 A(S')\times_{A(S)}A(S') \arrow[r,"\sim"] \arrow[d]
 & A(S')\times A(k[\operatorname{Ker}f]) \arrow[d]\\
 B(S')\times_{B(S)}B(S') \arrow[r,"\sim"']
 & B(S')\times B(k[\operatorname{Ker}f]).
\end{tikzcd}
\]
因此 $A(S')$、$B(S')$ 分别为 $A(S)$、$B(S)$ 上的主齐性集,
其结构群由
\[
 A(k[\operatorname{Ker}f])\xrightarrow{\ \sim\ }
 B(k[\operatorname{Ker}f]).
\]
认同. 若 $\varphi(S)$ 为双射, 则得 $\varphi(S')$ 为单射, 故为双射.
归纳完成. \hfill$\square$

\underline{推论 \SGARefTarget{sga7.sga7i.evi.n2-16}{2.16}.}
设 $\underline A$ 为 $\underline C_\Lambda$ 上的齐性群胚, 满足
\[
 \operatorname{Aut}_{\underline A(k[\varepsilon])}(1_\varepsilon)
\]
有限维. 对每个 $\xi\in\operatorname{ob}(\widehat{\underline A})$, 函子
\[
 \operatorname{Isom}_\xi:
 \widehat{\underline C}_{R\widehat\otimes_\Lambda R}
 \longrightarrow(\mathrm{Sets})
\]
前可表, 其中 $R=P(\xi)$.

\underline{证明.} 若 $\operatorname{Isom}_\xi$ 齐性,
则定理 \SGARefLink{sga7.sga7i.evi.n1-11}{1.11} 给出极小光滑函子
\[
 \Phi:\underline h_{R_1}\longrightarrow
 \underline{\operatorname{Isom}}_\xi.
\]
由引理 \SGARefLink{sga7.sga7i.evi.n2-15}{2.15}, $\Phi$ 是
$\widehat{\underline C}_{R\widehat\otimes_\Lambda R}$-等价,
故 $\operatorname{Isom}_\xi$ 前可表. 只须证明齐性.
考虑 $\widehat{\underline C}_{R\widehat\otimes_\Lambda R}$ 中的图
\[
\begin{tikzcd}[column sep=large,row sep=large]
 T_1 \arrow[dr] && T_2 \arrow[dl]\\
 & T &
\end{tikzcd}
\]
其中 $T_2\longrightarrow T$ 满射.
因 $\underline A$ 齐性, 纤维积
\[
 \xi_i(T_1)\times_{\xi_i(T)}\xi_i(T_2)
\]
在 $\underline A$ 中存在, 且有典范同构
\[
 \xi_i(T_1\times_TT_2)\xrightarrow{\ \sim\ }
 \xi_i(T_1)\times_{\xi_i(T)}\xi_i(T_2)
 \qquad(i=1,2)
\]
(参见 \SGARefLink{sga7.sga7i.evi.n2-6.item-a}{2.6}(a)). 从而
\[
 \operatorname{Isom}_\xi(T_1\times_TT_2)
 \longrightarrow
 \operatorname{Isom}_\xi(T_1)
 \times_{\operatorname{Isom}_\xi(T)}
 \operatorname{Isom}_\xi(T_2)
\]
为同构. \hfill$\square$

\underline{定理 \SGARefTarget{sga7.sga7i.evi.n2-17}{2.17}.}
设 $\underline A$ 为 $\underline C_\Lambda$ 上的余纤维群胚.
则它可由一个(正规化)光滑余群胚前表示当且仅当

%%%%%%%%%%  scan idx 79  |  folio 68  %%%%%%%%%%

(1) $\underline A$ 齐性;

(2) $[\underline A(k[\varepsilon])]$ 及
$\operatorname{Aut}_{\underline A(k[\varepsilon])}(1_\varepsilon)$
为 $k$ 上的有限维向量空间.

\underline{证明.} 首先假定 $\underline A$ 由光滑余群胚前表示.
即 $\underline A$ 在 $\underline C_\Lambda$-等价意义下等同于
\[
 \underline h=\underline h_{(R_0,R_1,\ldots)}
\]
其中 $(R_0,R_1,\ldots)$ 是 $\widehat{\underline C}_\Lambda$ 中的光滑余群胚.
由命题 \SGARefLink{sga7.sga7i.evi.n2-12}{2.12}(a), $\underline h$ 齐性.
而且,
\begin{align*}
 [\underline h(k[\varepsilon])]
 &=\operatorname{Coker}\!\left(
   \operatorname{Hom}_{\Lambda\text{-alg}}(R_1,k[\varepsilon])
   \mathrel{\substack{\xrightarrow{\ s_1\ }\\[-.6ex]
                       \xrightarrow[\ s_2\ ]{}}}
   \operatorname{Hom}_{\Lambda\text{-alg}}(R_0,k[\varepsilon])
   \right),\\[1ex]
 \operatorname{Aut}_{\underline h}(1_\varepsilon)
 &=\operatorname{Ker}\!\left(
   \operatorname{Hom}_{\Lambda\text{-alg}}(R_1,k[\varepsilon])
   \xrightarrow{\ s_1\ }
   \operatorname{Hom}_{\Lambda\text{-alg}}(R_0,k[\varepsilon])
   \right)\\
 &=\operatorname{Ker}\!\left(
   \operatorname{Hom}_{\Lambda\text{-alg}}(R_1,k[\varepsilon])
   \xrightarrow{\ s_2\ }
   \operatorname{Hom}_{\Lambda\text{-alg}}(R_0,k[\varepsilon])
   \right).
\end{align*}
二者均为 $k$ 上有限维向量空间, 因而 $\underline A$ 亦然.

反之, 假定 (1) 及 (2).
因 $\underline A$ 齐性且 $[\underline A(k[\varepsilon])]$ 有限维,
它有极小遍有形式对象 $\xi$. 置 $P(\xi)=R_0$, 并考虑
\[
 \operatorname{Isom}_\xi:
 \widehat{\underline C}_{R_0\widehat\otimes_\Lambda R_0}
 \longrightarrow(\mathrm{Sets}),
\]
由
\[
 T\longmapsto
 \operatorname{Isom}_{\underline A(T)}(\xi_1(T),\xi_2(T)),
\]
给出, 其中 $\xi_i(T)=\xi(f_i)$, $f_i:R_0\longrightarrow T$ 为图中复合映射
\[
\begin{tikzcd}[column sep=large]
 R_0 \arrow[r,shift left=2,"s_1"]
     \arrow[r,shift right=2,"s_2"']
 & R_0\widehat\otimes_\Lambda R_0 \arrow[r] & T.
\end{tikzcd}
\]
由推论 \SGARefLink{sga7.sga7i.evi.n2-16}{2.16}, 此函子由一个对象
\[
 R_1\in
 \operatorname{ob}(\widehat{\underline C}_{R_0\widehat\otimes_\Lambda R_0})
\]
及同构 $\xi_1(R_1)\xrightarrow{\varphi}\xi_2(R_1)$ 前表示, 满足
\[
 \Phi:\underline h_{R_1}\longrightarrow
 \underline{\operatorname{Isom}}_\xi
\]
为 $\widehat{\underline C}_{R_0\widehat\otimes_\Lambda R_0}$ 上的同构.
因而 $(R_0,R_1)$ 定义余群胚, 且
\[
 (\xi,\varphi):\underline h_{(R_0,R_1,\ldots)}
 \longrightarrow\underline A
\]
为 $\underline C_\Lambda$-等价.
因 $\xi$ 极小遍有, 复合
\[
 \underline h_{(R_0,R_0,\ldots)}
 \xrightarrow{\ \epsilon\ }
 \underline h_{(R_0,R_1,\ldots)}
 \xrightarrow{\ (\xi,\varphi)\ }
 \underline A
\]
极小光滑. 因 $(\xi,\varphi)$ 为等价, $\epsilon$ 也极小光滑.
故 $(R_0,R_1,\ldots)$ 是 $\widehat{\underline C}_\Lambda$ 中的正规化光滑余群胚.
\hfill$\square$

% END INCLUDED SGA7I_English_current_buildable_source_20260823/components/25_expose_VI_2_15_through_2_17.tex

% BEGIN INCLUDED SGA7I_English_current_buildable_source_20260823/components/26_expose_VI_2_18_through_2_22.tex
% R3 explicit revisions VI-A-05; see qa/FIDELITY-R3-CHANGES.json.
%%%%%%%%%%  scan idx 80  |  folio 69  %%%%%%%%%%

\underline{推论 \SGARefTarget{sga7.sga7i.evi.n2-18}{2.18}.}
函子 $F:\underline C_\Lambda\longrightarrow(\mathrm{Sets})$
前可表当且仅当它齐性且 $\dim F(k[\varepsilon])<\infty$.

\underline{推论 \SGARefTarget{sga7.sga7i.evi.n2-19}{2.19}.}
$\widehat{\underline C}_\Lambda$ 中的每个光滑余群胚都能在同构意义下正规化.

给定 $\widehat{\underline C}_\Lambda$ 中的余群胚 $(R_0,R_1,\ldots)$,
置 $\bar R_1=R_1/J$, 其中 $J$ 是 $R_1$ 中由
\[
 \{s_1(r)-s_2(r)\mid r\in R_0\}.
\]
生成的理想. 则 $(R_0,\bar R_1,\ldots)$ 带有诱导的合成法则,
定义 $R_0$ 上的形式群概形.

现设 $\underline A$ 为 $\underline C_\Lambda$ 上的余纤维群胚,
经由 $(\xi,\varphi)$ 由正规化光滑余群胚 $(R_0,R_1,\ldots)$ 前表示.
设
\[
 \underline C_k\lhook\joinrel\longrightarrow
 \underline C_{R_0}\lhook\joinrel\longrightarrow
 \underline C_{R_0\widehat\otimes_\Lambda R_0}
\]
为对应于满射
\[
 R_0\widehat\otimes_\Lambda R_0\longrightarrow R_0\longrightarrow k,
\]
的嵌入, 并以 $\operatorname{Aut}_\xi$ 和 $\operatorname{Aut}^{o}$
分别表示 $\operatorname{Isom}_\xi$ 在 $\underline C_{R_0}$ 和 $\underline C_k$ 上的限制.
于是 $\operatorname{Aut}_\xi(S)$ 为 $\xi(S)$ 的自同构群,
其中 $S\in\operatorname{ob}(\underline C_{R_0})$;
而 $\operatorname{Aut}^{o}(S)$ 为 $1_S$ 的自同构群,
其中 $S\in\operatorname{ob}(\underline C_k)$,
$1_S$ 为 $1\in\underline A(k)$ 在 $k\longrightarrow S$ 下的正像.

因 $R_1$ 在 $\underline C_{R_0\widehat\otimes_\Lambda R_0}$ 上前表示
$\operatorname{Isom}_\xi$, 故 $\bar R_1$ 在 $\underline C_{R_0}$ 上前表示
$\operatorname{Aut}_\xi$, 且 $k\otimes_{R_0}\bar R_1$ 在
$\underline C_k$ 上前表示 $\operatorname{Aut}^{o}$.
虽然 $\underline A$ 由光滑余群胚 $(R_0,R_1,\ldots)$ 前表示,
一般而言, $(R_0,\bar R_1,\ldots)$ 及
$(k,k\otimes_{R_0}\bar R_1,\ldots)$ 均不光滑. 更确切地:

\underline{推论 \SGARefTarget{sga7.sga7i.evi.n2-20}{2.20}.}
设 $\underline A$ 为 $\underline C_\Lambda$ 上的齐性余纤维群胚,
经由 $(\xi,\varphi)$ 由正规化光滑余群胚 $(R_0,R_1,\ldots)$ 前表示.

\begin{enumerate}[label=(\alph*)]
\item $(R_0,\bar R_1,\ldots)$ 及 $(k,k\otimes_{R_0}\bar R_1,\ldots)$
分别前表示 $\operatorname{Aut}_\xi$ 及 $\operatorname{Aut}^{o}$.

\item 下列断言等价:
\begin{enumerate}[label=(\roman*)]
\item 形式群概形 $\underline h_{(R_0,\bar R_1,\ldots)}$ 在 $R_0$ 上形式光滑;
\item 对 $\underline A$ 中每个满射箭头 $a'\longrightarrow a$, 映射
\[
 \operatorname{Aut}_{P(a')}(a')\longrightarrow
 \operatorname{Aut}_{P(a)}(a)
\]
满射;
\item $R_0$ 在 $\underline C_\Lambda$ 上前表示函子 $[\underline A]$.
\end{enumerate}
\end{enumerate}

%%%%%%%%%%  scan idx 81  |  folio 70  %%%%%%%%%%

\underline{证明.} (a) 已经指出.
对 (b), 等价性 (i)$\Leftrightarrow$(ii) 显然,
而 (ii)$\Leftrightarrow$(iii) 由 \SGARefLink{sga7.sga7i.evi.n2-7}{2.7} 及
\SGARefLink{sga7.sga7i.evi.n2-15}{2.15} 得到, 因为
\[
 \widetilde\xi:\underline h_{R_0}\longrightarrow[\underline A]
\]
极小光滑. \hfill$\square$

\underline{注 \SGARefTarget{sga7.sga7i.evi.n2-21}{2.21}.}
设 $\underline A$ 为 $\underline C_\Lambda$ 上的余纤维群胚,
$\underline A|\underline C_k$ 为
$p:\underline A\longrightarrow\underline C_\Lambda$ 沿
$\underline C_k\lhook\joinrel\longrightarrow\underline C_\Lambda$ 的拉回.
若 $\underline A$ 经由 $(\xi,\varphi)$ 由 $(R_0,R_1,\ldots)$ 前表示,
则 $\underline A|\underline C_k$ 经由
$(k\otimes_\Lambda\xi,k\otimes_\Lambda\varphi)$ 由
\[
 (k\otimes_\Lambda R_0,k\otimes_\Lambda R_1,\ldots)
\]
前表示. 此处 $k\otimes_\Lambda\xi$ 为 $\xi$ 在
$R_0\longrightarrow k\otimes_\Lambda R_0$ 下的正像,
而 $k\otimes_\Lambda\varphi$ 为 $\varphi$ 在映射
\[
 \operatorname{Isom}_{\underline A(R_1)}
   \bigl(\xi(s_1),\xi(s_2)\bigr)
 \longrightarrow
 \operatorname{Isom}_{\underline A(k\otimes_\Lambda R_1)}
   \bigl(k\otimes_\Lambda\xi(s_1),
         k\otimes_\Lambda\xi(s_2)\bigr).
\]
下的像. 若 $(R_0,R_1,\ldots)$ 正规化,
则 $(k\otimes_\Lambda R_0,k\otimes_\Lambda R_1,\ldots)$ 也正规化;
而且因 $R_1$ 在 $R_0$ 上光滑, $R_1$ 在 $R_0$ 上的相对维数等于
$k\otimes_\Lambda R_1$ 在 $k\otimes_\Lambda R_0$ 上的相对维数.

\underline{注 \SGARefTarget{sga7.sga7i.evi.n2-22}{2.22}.}
假定 $\underline A$ 分别由光滑余群胚 $(R_0,R_1,\ldots)$ 及
$(R'_0,R'_1,\ldots)$ 前表示, 其中 $(R_0,R_1,\ldots)$ 正规化.
得到同态
\[
 (i_0,i_1):(R_0,R_1,\ldots)\longrightarrow(R'_0,R'_1,\ldots).
\]
由 \SGARefLink{sga7.sga7i.evi.n2-10}{2.10}(b),
$R'_0$ 是 $R_0$ 上的形式幂级数环,
且 $(i_0,i_1)$ 诱导同构
\[
 R'_0\otimes_{R_0}(R_0,R_1,\ldots)
 \xrightarrow{\ \sim\ }(R'_0,R'_1,\ldots).
\]

% END INCLUDED SGA7I_English_current_buildable_source_20260823/components/26_expose_VI_2_18_through_2_22.tex
% BEGIN INCLUDED SGA7I_English_current_buildable_source_20260823/components/27_expose_VI_section_3_opening_through_3_1_b.tex
%%%%%%%%%%  scan idx 81  |  folio 70  %%%%%%%%%%

\section*{\SGARefTarget{sga7.sga7i.evi.n3}{3}. 凝聚层 $\underline D_X^i$}

形变函子是一个重要的函子, 它通常并非前可表, 但仍有遍有形式对象.
为此回顾交换环扩张及其障碍理论的一些基本事实.
本节除系统使用景上的 Grothendieck 上同调之外, 或许没有新的内容.
其目的是使本报告尽可能自足; 实际上我们所证明的内容略多于后文所需.
基本结果陈述于定理 \SGARefLink{sga7.sga7i.evi.n3-12}{3.12}.
全文中所有环均为有单位元的交换环.

%%%%%%%%%%  scan idx 82  |  folio 71  %%%%%%%%%%

给定环 $R$, 用 $\underline C_R$ 表示 $R$-代数范畴.
对每个 $A\in\operatorname{Ob}(\underline C_R)$, 置
\[
 \operatorname{Cov}(A)=
 \{\text{满射箭头 }\varphi:A'\longrightarrow A
   \text{ 位于 }\underline C_R\}.
\]
这在 $\underline C_R$ 上定义预拓扑, 因而定义一个景(site).
对每个 $A\in\operatorname{Ob}(\underline C_R)$,
设 $\underline C_{A|R}$ 为 $\underline C_R$ 中靶为 $A$ 的箭头范畴.
赋予 $\underline C_{A|R}$ 由遗忘函子
\[
 j_A:\underline C_{A|R}\longrightarrow\underline C_R.
\]
诱导的拓扑. 范畴 $\underline C_{A|R}$ 有终对象
$(1_A:A\longrightarrow A)$ 及始对象 $(i_A:R\longrightarrow A)$,
且容许纤维积与余积.
无混淆之虞时, 将 $\underline C_{A|R}$ 的对象 $(B\longrightarrow A)$ 简记为 $B$.

设 $\mathscr C$ 为阿贝尔范畴.
按定义, $\underline C_{A|R}$ 上的 $\mathscr C$-值层是函子
\[
 F:\underline C_{A|R}^{\mathrm{op}}\longrightarrow\mathscr C
\]
使对 $\underline C_{A|R}$ 中每个满射箭头 $B'\longrightarrow B$, 序列
\[
 0\longrightarrow F(B)\longrightarrow F(B')
 \mathrel{\substack{\longrightarrow\\[-.6ex]\longrightarrow}}
 F(B'\mathbin{\times_B}B')
\]
正合. \underline{加性}预层是保持余积的预层; 因此
\[
 F(B_1\otimes_R B_2)=F(B_1)\oplus F(B_2).
\]
例如, $A$-模 $M$ 定义函子
\[
 \operatorname{Der}_R(-,M):
 \underline C_{A|R}^{\mathrm{op}}\longrightarrow(A\text{-mod}),
\]
将 $B$ 送到 $R$-线性导子 $\operatorname{Der}_R(B,M)$,
其中经由结构映射 $B\longrightarrow A$ 将 $M$ 看作 $B$-模.
这显然是 $\underline C_{A|R}$ 上的加性层.
此处 $(A\text{-mod})$ 表示 $A$-模范畴.

对 $\underline C_{A|R}$ 上取值于 $(A\text{-mod})$ 的预层 $F$,
若只要 $B$ 是有限型 $R$-代数, $F(B)$ 就是有限型 $A$-模,
则称其为\underline{有限型}.
例如, 若 $A$ 诺特且 $M$ 是有限型 $A$-模,
则 $\operatorname{Der}_R(-,M)$ 是 $\underline C_{A|R}$ 上的有限型加性层.

再记录一些记号及同调代数的基本事实.

\underline{\SGARefTarget{sga7.sga7i.evi.n3-1.item-a}{\SGARefAliasTarget{sga7.sga7i.evi.n3-1}3.1}(a).}
本节余下部分, 分别以 $\widetilde{\underline C}_{A|R}$ 及
$\overset{\vee}{\underline C}_{A|R}$ 表示取值于 $(A\text{-mod})$ 的层与预层范畴.
有左正合的包含函子
\[
 i_{A|R}:\widetilde{\underline C}_{A|R}
 \longrightarrow\overset{\vee}{\underline C}_{A|R}
\]
及正合的伴随层函子
\[
 G_{A|R}:\overset{\vee}{\underline C}_{A|R}
 \longrightarrow\widetilde{\underline C}_{A|R},
\]
它是 $i_{A|R}$ 的左伴随.

%%%%%%%%%%  scan idx 83  |  folio 72  %%%%%%%%%%

对象 $B\in\operatorname{Ob}(\underline C_{A|R})$ 诱导连续函子
$j:\underline C_{B|R}\longrightarrow\underline C_{A|R}$, 因而得到交换图
\begin{equation}\tag{*}
\begin{tikzcd}[column sep=huge,row sep=huge]
 \widetilde{\underline C}_{A|R}
   \arrow[r,"\widetilde j"] \arrow[d,"i_{A|R}"']
 & \widetilde{\underline C}_{B|R} \arrow[d] \\
 \overset{\vee}{\underline C}_{A|R}
   \arrow[r,"j^{\vee}"']
 & \overset{\vee}{\underline C}_{B|R}.
\end{tikzcd}
\end{equation}
限制函子 $\widetilde j$ 及 $j^{\vee}$ 均正合,
因为景 $\underline C_{B|R}$ 就是 $\underline C_{A|R}$
中靶为 $B$ 的箭头范畴, 赋予诱导拓扑.

环同态 $S\longrightarrow R$ 诱导连续函子
$\beta:\underline C_{A|R}\longrightarrow\underline C_{A|S}$,
因而得到交换图
\begin{equation}\tag{**}
\begin{tikzcd}[column sep=huge,row sep=huge]
 \widetilde{\underline C}_{A|S}
   \arrow[r,"\widetilde\beta"] \arrow[d,"i_{A|S}"']
 & \widetilde{\underline C}_{A|R} \arrow[d] \\
 \overset{\vee}{\underline C}_{A|S}
   \arrow[r,"\beta^{\vee}"']
 & \overset{\vee}{\underline C}_{A|R}.
\end{tikzcd}
\end{equation}
注意 $\widetilde\beta$ 未必正合.

\underline{\SGARefTarget{sga7.sga7i.evi.n3-1.item-b}{3.1}(b).}
将包含函子
\[
 i_{A|R}:\widetilde{\underline C}_{A|R}
 \longrightarrow\overset{\vee}{\underline C}_{A|R}
\]
的右导出函子记为 $\underline H^{\bullet}_{A|R}$.
对 $\underline C_{A|R}$ 上每个层 $F$, 置
\[
 H^{\bullet}(A|R,F)=
 [\underline H^{\bullet}_{A|R}(F)](A).
\]

对 $\operatorname{Cov}(B)$ 中的 $B'\longrightarrow B$, 定义
\[
 H^0(\{B'\longrightarrow B\},F)=
 \ker\!\left(
 F(B')\mathrel{\substack{\longrightarrow\\[-.6ex]\longrightarrow}}
 F(B'\mathbin{\times_B}B')
 \right).
\]
将 $H^0(\{B'\longrightarrow B\},-)$ 的右导出函子记为
$H^{\bullet}(\{B'\longrightarrow B\},-)$.
这些就是下述 \v{C}ech 半单纯模的上同调群

%%%%%%%%%%  scan idx 84  |  folio 73  %%%%%%%%%%

\[
\begin{tikzcd}[column sep=large]
 F(B')
   \arrow[r,shift left=2] \arrow[r,shift right=2]
 & F(B'\times_BB')
   \arrow[r,shift left=3] \arrow[r] \arrow[r,shift right=3]
 & F(B'\times_BB'\times_BB')
   \arrow[r,shift left=4.5] \arrow[r,shift left=1.5]
   \arrow[r,shift right=1.5] \arrow[r,shift right=4.5]
 & \cdots .
\end{tikzcd}
\]

将
\[
 \check{\underline H}^{,0}_{A|R}:
 \overset{\vee}{\underline C}_{A|R}
 \longrightarrow\overset{\vee}{\underline C}_{A|R}
\]
的右导出函子记为 $\check{\underline H}^{\,\bullet}_{A|R}$, 其中
\[
 [\check{\underline H}^{,0}_{A|R}(F)](B)=
 \varinjlim_{\{B'\to B\}\in\operatorname{Cov}(B)}
 H^0(\{B'\longrightarrow B\},F).
\]
对 $\underline C_{A|R}$ 上的每个预层 $F$, 置
\[
 \check H^{\bullet}(A|R,F)=
 [\check{\underline H}^{\,\bullet}_{A|R}(F)](A)=
 \varinjlim_{\{A'\to A\}\in\operatorname{Cov}(A)}
 H^{\bullet}(\{A'\longrightarrow A\},F).
\]
对 $\underline C_{A|R}$ 上的每个层 $F$, 有通常的谱序列
\[
 \check H^p\bigl(A|R,\underline H^q_{A|R}(F)\bigr)
 \Longrightarrow H^{p+q}(A|R,F).
\]
此处对每个 $q>0$ 均有 $\check H^0(A|R,\underline H^q_{A|R}(F))=0$.
从而,
\[
 \check H^1(A|R,F)=H^1(A|R,F),
\]
且序列
\[
 0\longrightarrow\check H^2(A|R,F)
 \longrightarrow H^2(A|R,F)
 \longrightarrow
 \check H^1\bigl(A|R,\underline H^1_{A|R}(F)\bigr)
 \longrightarrow\check H^3(A|R,F)
\]
正合.

% END INCLUDED SGA7I_English_current_buildable_source_20260823/components/27_expose_VI_section_3_opening_through_3_1_b.tex

% BEGIN INCLUDED SGA7I_English_current_buildable_source_20260823/components/28_expose_VI_3_1_c_through_3_4.tex
\underline{\SGARefTarget{sga7.sga7i.evi.n3-1.item-c}{3.1}(c).}
若 $B$ 为 $\underline C_{A|R}$ 的对象,
则 \SGARefLink{sga7.sga7i.evi.n3-1.item-a}{3.1}(a) 中的图 $(*)$,
连同 $\widetilde j$ 及 $j^{\vee}$ 的正合性, 给出典范同构
\[
 H^{\bullet}(B|R,\widetilde jF)=
 [\underline H^{\bullet}_{A|R}(F)](B).
\]
另一方面, 环同态 $S\longrightarrow R$ 诱导
\SGARefLink{sga7.sga7i.evi.n3-1.item-a}{3.1}(a) 中的图 $(**)$.
函子 $\widetilde\beta$ 未必正合; 但 $\beta$ 保持纤维积, 故得到谱序列
\[
 \underline H^p_{A|R}(R^q\widetilde\beta F)
 \Longrightarrow
 \beta^{\vee}\underline H^{p+q}_{A|S}(F),
\]
或在 $A$ 处取值, 得到
\[
 H^p(A|R,R^q\widetilde\beta F)
 \Longrightarrow H^{p+q}(A|S,F).
\]
因有典范同构
\[
 \widetilde\beta=
 G_{A|R}\circ\beta^{\vee}\circ i_{A|S}
\]
且函子 $\beta^{\#}=G_{A|R}\circ\beta^{\vee}$ 正合, 有

%%%%%%%%%%  scan idx 85  |  folio 74  %%%%%%%%%%

\[
 R^{\bullet}\widetilde\beta F=
 \beta^{\#}\underline H^{\bullet}_{A|S}(F).
\]
因此上述谱序列可写为
\[
 H^p\bigl(A|R,\beta^{\#}\underline H^q_{A|S}(F)\bigr)
 \Longrightarrow H^{p+q}(A|S,F).
\]

\underline{引理 \SGARefTarget{sga7.sga7i.evi.n3-2}{3.2}.}
设 $B\in\operatorname{Ob}(\underline C_{A|R})$,
并在 $\underline C_{B|R}$ 中取一个满射箭头 $P\longrightarrow B$,
其中 $P$ 是 $R$ 上的多项式代数.
则对每个 $F\in\operatorname{Ob}(\overset{\vee}{\underline C}_{A|R})$, 自然映射
\[
 H^{\bullet}(\{P\longrightarrow B\},F)
 \longrightarrow[\check{\underline H}^{\,\bullet}(F)](B)
\]
为同构.

\underline{证明.} 只须证明自然映射
\[
 H^0(\{P\longrightarrow B\},F)
 \longrightarrow[\check{\underline H}^{,0}(F)](B)
 =\varinjlim_{\{B'\to B\}\in\operatorname{Cov}(B)}
   H^0(\{B'\longrightarrow B\},F)
\]
为同构. 这是显然的, 因为对每个
$\{B'\longrightarrow B\}\in\operatorname{Cov}(B)$, 存在覆盖间的映射
\[
 \{P\longrightarrow B\}\longrightarrow
 \{B'\longrightarrow B\}.
\]
\hfill$\square$

\underline{推论 \SGARefTarget{sga7.sga7i.evi.n3-3}{3.3}.}

\begin{enumerate}[label=(\alph*)]
\item 假定 $A$ 诺特. 若
$F\in\operatorname{Ob}(\widetilde{\underline C}_{A|R})$ 有限型,
则对每个 $n$, $\underline H^n(F)$ 均有限型.

\item 若 $P$ 是 $R$ 上的多项式代数, 则
\[
 H^n(P|R,F)=0
\]
对每个 $n>0$ 及每个 $F\in\operatorname{Ob}(\widetilde{\underline C}_{P|R})$ 成立.

\item $\widetilde{\underline C}_{A|R}$ 中的序列
\[
 0\longrightarrow F'\longrightarrow F\longrightarrow F''
 \longrightarrow0
\]
正合, 当且仅当对每个多项式 $R$-代数
$P\in\operatorname{Ob}(\underline C_{A|R})$, 序列
\[
 0\longrightarrow F'(P)\longrightarrow F(P)
 \longrightarrow F''(P)\longrightarrow0
\]
正合.
\end{enumerate}

\underline{证明.}

(a) 由谱序列
\[
 \check H^p\bigl(B|R,\underline H^q(F)\bigr)
 \Longrightarrow H^{p+q}(B|R,F)
\]
及 $q>0$ 时的等式 $\check H^0(\underline H^q(F))=0$,
只须证明对每个 $n$ 及每个有限型预层 $G$,
$\check{\underline H}^{,n}(G)$ 均有限型.
由引理 \SGARefLink{sga7.sga7i.evi.n3-2}{3.2},
\[
 \check H^{\bullet}(B|R,G)=
 H^{\bullet}(\{P\longrightarrow B\},G),
\]
其中 $P$ 为 $R$ 上的多项式代数,
$\{P\longrightarrow B\}\in\operatorname{Cov}(B)$.
后者是下述复形的上同调:

%%%%%%%%%%  scan idx 86  |  folio 75  %%%%%%%%%%

\[
 \cdots
 G(P\times_BP\times_BP)
 \mathrel{\substack{\longleftarrow\\[-.6ex]\longleftarrow\\[-.6ex]\longleftarrow}}
 G(P\times_BP)
 \mathrel{\substack{\longleftarrow\\[-.6ex]\longleftarrow}}
 G(P).
\]
若 $B$ 为有限型 $R$-代数, 可取 $P$ 为 $R$ 上有限型的.
所有迭代纤维积 $P\times_B\cdots\times_BP$ 于是都是有限型 $R$-代数.
若 $G$ 有限型, 则它们在 $G$ 下的像为有限型 $A$-模.
故因 $A$ 诺特, $H^{\bullet}(\{P\longrightarrow B\},G)$ 是有限型 $A$-模.

(b) 设 $P$ 为 $R$ 上的多项式代数.
每个满射 $R$-代数映射 $B\longrightarrow P$ 有截面, 故
\[
 H^i(\{B\longrightarrow P\},G)=0\qquad(i>0).
\]
由此, 对每个 $i>0$ 及 $\underline C_{P|R}$ 上的每个预层 $G$,
均有 $\check H^i(P|R,G)=0$. 从而谱序列
\[
 \check H^p\bigl(P|R,\underline H^q_{P|R}(F)\bigr)
 \Longrightarrow H^{p+q}(P|R,F)
\]
退化, 且对每个 $n$, 有
\[
 \check H^0\bigl(P|R,\underline H^n(F)\bigr)
 \simeq H^n(P|R,F)
\]
然而, 对层 $F$, 只要 $n>0$ 就有
$\check H^0(P|R,\underline H^n(F))=0$.
因此对每个 $n>0$, $H^n(P|R,F)=0$.

(c) 若 $0\longrightarrow F'\longrightarrow F\longrightarrow F''\longrightarrow0$
在 $\widetilde{\underline C}_{A|R}$ 中正合, 则序列
\[
 0\longrightarrow F'(B)\longrightarrow F(B)
 \longrightarrow F''(B)
 \longrightarrow H^1(B|R,F')\longrightarrow H^1(B|R,F)
 \longrightarrow\cdots
\]
正合. 若 $P\in\operatorname{Ob}(\underline C_{A|R})$ 为 $R$ 上的多项式代数,
则由 (b), $H^1(P|R,F')=0$, 故
\[
 0\longrightarrow F'(P)\longrightarrow F(P)
 \longrightarrow F''(P)\longrightarrow0
\]
正合.

反之, 设 $\underline P_{A|R}$ 为 $\underline C_{A|R}$ 的满子范畴, 其对象为
\[
 \operatorname{Ob}(\underline P_{A|R})=
 \{P\in\operatorname{Ob}(\underline C_{A|R})\mid
   P\text{ 是多项式代数, 基环为 }R\},
\]
赋予诱导拓扑, 并设 $p:\underline P_{A|R}\longrightarrow\underline C_{A|R}$ 为包含函子.
$\underline C_{A|R}$ 的每个对象都被 $\underline P_{A|R}$ 的对象覆盖;
因此比较引理表明
\[
 \widetilde p:\widetilde{\underline C}_{A|R}
 \longrightarrow\widetilde{\underline P}_{A|R}
\]
为范畴等价. 若将 $F'\longrightarrow F\longrightarrow F''$
在每个 $P\in\operatorname{Ob}(\underline P_{A|R})$ 处取值后得到的序列均正合, 则
\[
 0\longrightarrow\widetilde pF'\longrightarrow\widetilde pF
 \longrightarrow\widetilde pF''\longrightarrow0
\]
在 $\widetilde{\underline P}_{A|R}$ 中正合,
因而原序列在 $\widetilde{\underline C}_{A|R}$ 中正合. \hfill$\square$

%%%%%%%%%%  scan idx 87  |  folio 76  %%%%%%%%%%

\underline{注 \SGARefTarget{sga7.sga7i.evi.n3-4}{3.4}.}
$\underline P_{A|R}$ 中的每个满射箭头均有截面.
立即得到 $\underline P_{A|R}$ 上的每个预层都是层; 即
\[
 \widetilde{\underline P}_{A|R}=
 \overset{\vee}{\underline P}_{A|R}.
\]
从而限制函子
\[
 \widetilde{\underline C}_{A|R}
 \longrightarrow\overset{\vee}{\underline P}_{A|R}
\]
为范畴等价.

% END INCLUDED SGA7I_English_current_buildable_source_20260823/components/28_expose_VI_3_1_c_through_3_4.tex
% BEGIN INCLUDED SGA7I_English_current_buildable_source_20260823/components/29_expose_VI_3_5_through_3_6.tex
% R3 explicit revisions VI-A-03; see qa/FIDELITY-R3-CHANGES.json.
\underline{引理 \SGARefTarget{sga7.sga7i.evi.n3-5}{3.5}.}
考虑 $R$-代数所成的图
\[
\begin{tikzcd}[column sep=large,row sep=large]
 & C \arrow[d]\\
 B' \arrow[r,"\text{满射}"'] & B
\end{tikzcd}
\]
并设 $S$ 为 $R$-代数.

\begin{enumerate}[label=(\roman*)]
\item 若 $\operatorname{Tor}_1^R(B,S)=0$, 则典范映射
\[
 (B'\times_BC)\otimes_RS
 \longrightarrow
 (B'\otimes_RS)\times_{B\otimes_RS}(C\otimes_RS)
\]
为同构.

\item 假定 $B'$ 在 $R$ 上平坦. 若
\[
 \operatorname{Tor}_i^R(C,S)=0\quad(0<i<n)
 \qquad\text{且}\qquad
 \operatorname{Tor}_i^R(B,S)=0\quad(0<i\leq n),
\]
则
\[
 \operatorname{Tor}_i^R(B'\times_BC,S)=0\qquad(0<i<n).
\]
\end{enumerate}

\underline{证明.} 设 $0\longrightarrow J\longrightarrow B'\longrightarrow B\longrightarrow0$ 正合.

(i) 若 $\operatorname{Tor}_1^R(B,S)=0$, 则
\[
 0\longrightarrow J\otimes_RS\longrightarrow B'\otimes_RS
 \longrightarrow B\otimes_RS\longrightarrow0
\]
正合, 从而
\[
 0\longrightarrow J\otimes_RS
 \longrightarrow
 (B'\otimes_RS)\times_{B\otimes_RS}(C\otimes_RS)
 \longrightarrow C\otimes_RS\longrightarrow0.
\]
也正合. 因此得到行正合的交换图:
\begin{center}
\scriptsize
\begin{tikzcd}[column sep=small,row sep=large]
 & J\otimes_RS \arrow[r] \arrow[d,equal]
 & (B'\times_BC)\otimes_RS \arrow[r] \arrow[d]
 & C\otimes_RS \arrow[r] \arrow[d,equal]
 & 0\\
 0 \arrow[r] & J\otimes_RS \arrow[r]
 & (B'\otimes_RS)\times_{B\otimes_RS}(C\otimes_RS) \arrow[r]
 & C\otimes_RS \arrow[r] & 0.
\end{tikzcd}
\end{center}

%%%%%%%%%%  scan idx 88  |  folio 77  %%%%%%%%%%

由此得到
\[
 (B'\times_BC)\otimes_RS
 \longrightarrow
 (B'\otimes_RS)\times_{B\otimes_RS}(C\otimes_RS)
\]
为同构.

(ii) 因 $B'$ 在 $R$ 上平坦, 且对 $0<i\leq n$,
$\operatorname{Tor}_i^R(B,S)=0$, Tor 长正合序列给出
\[
 \operatorname{Tor}_i^R(J,S)=0\qquad(0<i<n).
\]
对
\[
 0\longrightarrow J\longrightarrow B'\times_BC
 \longrightarrow C\longrightarrow0
\]
应用同一长正合序列, 并使用对 $C$ 的假设, 得到
\[
 \operatorname{Tor}_i^R(B'\times_BC,S)=0\qquad(0<i<n).
\]
\hfill$\square$

\underline{推论 \SGARefTarget{sga7.sga7i.evi.n3-6}{3.6}.}

\begin{enumerate}[label=(\alph*)]
\item 设 $A$ 和 $S$ 为 $R$-代数.
若 $\operatorname{Tor}_1^R(A,S)=0$, 则对每个
$G\in\operatorname{Ob}(\overset{\vee}{\underline C}_{A\otimes_RS|S})$, 典范映射
\[
 \check H^{\bullet}(A\otimes_RS|S,G)
 \longrightarrow\check H^{\bullet}(A|R,\check SG)
\]
为同构. 此处
\[
 S:\underline C_{A|R}\longrightarrow
 \underline C_{A\otimes_RS|S},
 \qquad X\longmapsto X\otimes_RS.
\]
特别地, 若 $S$ 在 $R$ 上平坦, 则
\[
 \check S\,\check{\underline H}^{\,\bullet}(G)
 \longrightarrow
 \check{\underline H}^{\,\bullet}(\check SG)
\]
对每个这样的 $G$ 都是同构.

\item 设 $P\longrightarrow B$ 为 $R$-代数的满射映射,
其中 $P$ 为 $R$ 上的多项式代数. 若
\[
 \operatorname{Tor}_i^R(B,S)=0\qquad(0<i\leq n),
\]
则对任意有限个 $P$ 在 $B$ 上的迭代纤维积,
\[
 \operatorname{Tor}_i^R
 (P\times_B P\times_B\cdots\times_B P,S)=0
 \qquad(0<i<n).
\]
\end{enumerate}

\underline{证明.}

(a) 给定 $B\in\operatorname{Ob}(\underline C_{A|R})$,
取覆盖 $\{P\longrightarrow B\}\in\operatorname{Cov}(B)$,
其中 $P$ 为 $R$ 上的多项式代数.
若 $\operatorname{Tor}_1^R(B,S)=0$, 则由引理 \SGARefLink{sga7.sga7i.evi.n3-5}{3.5}(i),
对每个 $n$, 自然映射
\[
 \left(\overbrace{P\times_B\cdots\times_BP}^{n}\right)\otimes_RS
 \longrightarrow
 \overbrace{(P\otimes_RS)\times_{B\otimes_RS}\cdots
 \times_{B\otimes_RS}(P\otimes_RS)}^{n}
\]
为同构. 从而,
\begin{align*}
 \check H^{\bullet}(B|R,\check SG)
 &=H^{\bullet}(\{P\longrightarrow B\},\check SG)\\
 &=H^{\bullet}(\{P\otimes_RS\longrightarrow B\otimes_RS\},G)\\
 &=\check H^{\bullet}(B\otimes_RS|S,G).
\end{align*}
若 $S$ 在 $R$ 上平坦, 则对每个
$X\in\operatorname{Ob}(\underline C_{A|R})$,
$\operatorname{Tor}_1^R(X,S)=0$.
上述同构因此关于 $X$ 自然, 给出
\[
 \check S\,\check{\underline H}^{\,\bullet}(G)
 \xrightarrow{\ \sim\ }
 \check{\underline H}^{\,\bullet}(\check SG).
\]

(b) 置
\[
 P^{(m)}=\overbrace{P\times_B\cdots\times_BP}^{m+1}.
\]
$m=0$ 时无须证明, 因 $P$ 在 $R$ 上平坦.
归纳假定对 $0<i<n$ 有 $\operatorname{Tor}_i^R(P^{(m-1)},S)=0$.
因 $P\longrightarrow B$ 满射且 $P^{(m)}=P\times_BP^{(m-1)}$,
由引理 \SGARefLink{sga7.sga7i.evi.n3-5}{3.5}(ii),
\[
 \operatorname{Tor}_i^R(P^{(m)},S)=0\qquad(0<i<n).
\]
\hfill$\square$

% END INCLUDED SGA7I_English_current_buildable_source_20260823/components/29_expose_VI_3_5_through_3_6.tex
% BEGIN INCLUDED SGA7I_English_current_buildable_source_20260823/components/30_expose_VI_proposition_3_7.tex
%%%%%%%%%%  scan idx 89  |  folio 78  %%%%%%%%%%

\underline{命题 \SGARefTarget{sga7.sga7i.evi.n3-7}{3.7}.}

\begin{enumerate}[label=(\alph*)]
\item 设 $A$ 和 $S$ 为 $R$-代数. 若
\[
 \operatorname{Tor}_i^R(A,S)=0\qquad(0<i\leq n),
\]
则对 $\underline C_{A\otimes_RS|S}$ 上每个层 $F$, 典范映射
\[
 H^i(A\otimes_RS|S,F)
 \longrightarrow H^i(A|R,\widetilde SF)
\]
对每个 $i\leq n$ 均为同构. 此处
\[
 S:\underline C_{A|R}\longrightarrow
 \underline C_{A\otimes_RS|S},
 \qquad X\longmapsto X\otimes_RS.
\]

\item 设 $P$ 为 $R$ 上的多项式代数.
则对 $\underline C_{A\otimes_RP|R}$ 上每个加性层 $F$, 典范映射
\[
 H^n(A\otimes_RP|R,F)\longrightarrow H^n(A|R,F)
\]
对每个 $n>0$ 均为同构.

\item 设 $S\longrightarrow R\longrightarrow A$ 为环同态.
则对 $\underline C_{A|S}$ 上每个加性层 $F$, 有正合序列
\begin{multline*}
0\longrightarrow H^0(A|R,F_R)\longrightarrow H^0(A|S,F)
\longrightarrow H^0(R|S,F)\longrightarrow H^1(A|R,F_R)\\
\longrightarrow H^1(A|S,F)\longrightarrow H^1(R|S,F)
\longrightarrow\cdots\longrightarrow H^n(A|R,F_R)\\
\longrightarrow H^n(A|S,F)\longrightarrow H^n(R|S,F)
\longrightarrow H^{n+1}(A|R,F_R)\longrightarrow\cdots,
\end{multline*}
其中 $F_R$ 是 $\underline C_{A|R}$ 上的层,
经由遗忘函子 $\underline C_{A|R}\longrightarrow\underline C_{A|S}$,
对每个 $X\in\operatorname{Ob}(\underline C_{A|R})$ 定义为
\[
 F_R(X)=\ker\bigl(F(X)\longrightarrow F(R)\bigr)
\]
\end{enumerate}

\underline{证明.}

(a) 考虑连续函子
\[
 S:\underline C_{A|R}\longrightarrow
 \underline C_{A\otimes_RS|S},
 \qquad S(X)=X\otimes_RS.
\]
因 $S$ 将 $R$ 上的多项式代数送至 $S$ 上的多项式代数,
由 \SGARefLink{sga7.sga7i.evi.n3-3}{3.3}(c), 函子
\[
 \widetilde S:\widetilde{\underline C}_{A\otimes_RS|S}
 \longrightarrow\widetilde{\underline C}_{A|R}
\]
正合. 因而得到典范映射
\[
 \check S\,\underline H^{\bullet}_{A\otimes_RS|S}(F)
 \longrightarrow
 \underline H^{\bullet}_{A|R}(\widetilde SF)
\]
及下列谱序列的交换图:
\begin{center}
\small
\begin{tikzcd}[column sep=huge,row sep=large]
E_S^{p,q}(S(X))=
 \check H^p\bigl(X\otimes_RS|S,
 \underline H^q_{A\otimes_RS|S}(F)\bigr)
 \arrow[r,Rightarrow] \arrow[d]
& H^{p+q}(X\otimes_RS|S,F) \arrow[dd]\\
\check H^p\bigl(X|R,
 \check S\,\underline H^q_{A\otimes_RS|S}(F)\bigr)
 \arrow[d] & \\
E_R^{p,q}(X)=
 \check H^p\bigl(X|R,
 \underline H^q_{A|R}(\widetilde SF)\bigr)
 \arrow[r,Rightarrow]
& H^{p+q}(X|R,\widetilde SF).
\end{tikzcd}
\end{center}

%%%%%%%%%%  scan idx 90  |  folio 79  %%%%%%%%%%

由 \SGARefLink{sga7.sga7i.evi.n3-6}{3.6}(a), 只要
$\operatorname{Tor}_1^R(X,S)=0$, 映射
\[
 E_S^{p,q}(S(X))\longrightarrow
 \check H^p\bigl(X|R,
 \check S\,\underline H^q_{A\otimes_RS|S}(F)\bigr)
\]
就是同构. 归纳假定对每个 $q<n$ 及每个满足
\[
 \operatorname{Tor}_i^R(X,S)=0\qquad(0<i\leq n).
\]
的 $X\in\operatorname{Ob}(\underline C_{A|R})$, 映射
\[
 [\check S\,\underline H^q_{A\otimes_RS|S}(F)](X)
 \longrightarrow
 [\underline H^q_{A|R}(\widetilde SF)](X)
\]
均为同构. 由 \SGARefLink{sga7.sga7i.evi.n3-6}{3.6}(b),
在同样的 Tor 假设下, 对每个 $q<n$ 及每个 $p$,
$E_S^{p,q}(S(X))\longrightarrow E_R^{p,q}(X)$ 均为同构.
而且 $E_S^{0,n}(S(X))\longrightarrow E_R^{0,n}(X)$ 为同构,
因为两端均为零. 故只要对 $0<i\leq n$ 有 $\operatorname{Tor}_i^R(X,S)=0$,
\[
 H^n(X\otimes_RS|S,F)
 \longrightarrow H^n(X|R,\widetilde SF)
\]
就是同构. 等价地, 在此假设下,
\[
 [\check S\,\underline H^i_{A\otimes_RS|S}(F)](X)
 \longrightarrow
 [\underline H^i_{A|R}(\widetilde SF)](X)
\]
对每个 $i\leq n$ 均为同构.

(b) 设
\[
 f:\underline C_{A|R}\longrightarrow
 \underline C_{A\otimes_RP|R},
 \qquad f(X)=X\otimes_RP.
\]
在 $\underline C_{A|R}$ 中取满射箭头 $P'\longrightarrow X$,
其中 $P'$ 为 $R$ 上的多项式代数. 则
\[
 P'\otimes_RP\longrightarrow X\otimes_RP
\]
为 $\underline C_{A\otimes_RP|R}$ 中的满射箭头,
且 $P'\otimes_RP$ 仍为 $R$ 上的多项式代数.
因 $P$ 在 $R$ 上平坦, 由 \SGARefLink{sga7.sga7i.evi.n3-5}{3.5}(i),
\begin{align*}
 \check H^{\bullet}(X|R,\check fG)
 &=H^{\bullet}(\{P'\longrightarrow X\},\check fG)\\
 &=H^{\bullet}(\{P'\otimes_RP\longrightarrow X\otimes_RP\},G)\\
 &=\check H^{\bullet}(X\otimes_RP|R,G)
\end{align*}
对每个 $G\in\operatorname{Ob}(\overset{\vee}{\underline C}_{A\otimes_RP|R})$ 成立.
故 $\widetilde f$ 将无上同调层送至无上同调层.
另一方面, $f$ 将 $R$ 上的多项式代数送至 $R$ 上的多项式代数, 故由
\SGARefLink{sga7.sga7i.evi.n3-3}{3.3}(c),
\[
 \widetilde f:\widetilde{\underline C}_{A\otimes_RP|R}
 \longrightarrow\widetilde{\underline C}_{A|R}
\]
正合. 由此对每个 $F\in\operatorname{Ob}(\widetilde{\underline C}_{A\otimes_RP|R})$,
\[
 H^{\bullet}(A\otimes_RP|R,F)
 \longrightarrow H^{\bullet}(A|R,\widetilde fF)
\]
为同构.

对每个 $X\in\operatorname{Ob}(\underline C_{A|R})$, 有典范映射
\[
 (\widetilde fF)(X)=F(X\otimes_RP)
 \longrightarrow F(X)\oplus F(P),
\]
因而有典范映射 $\widetilde fF\longrightarrow F\oplus F(P)$,
其中 $F(P)$ 表示常层, 右侧的 $F$ 表示 $F$ 在 $\underline C_{A|R}$ 上的限制.
若 $F$ 加性, 此映射为同构. 因此对每个 $n>0$,
\[
 H^n(A|R,\widetilde fF)
 \xrightarrow{\ \sim\ }
 H^n(A|R,F\oplus F(P))=H^n(A|R,F),
\]
因为常层 $F(P)$ 为松层. 这证明了 (b).

%%%%%%%%%%  scan idx 91  |  folio 80  %%%%%%%%%%

(c) 此论证由 Rinehart 告知.
考虑遗忘函子 $\beta:\underline C_{A|R}\longrightarrow\underline C_{A|S}$.
有谱序列
\[
 E^{p,q}=H^p\bigl(A|R,
 \beta^{\#}\underline H^q_{A|S}(F)\bigr)
 \Longrightarrow H^{p+q}(A|S,F),
\]
其中
\[
 \beta^{\#}=G_{A|R}\circ\beta^{\vee},
 \qquad
 G_{A|R}:\overset{\vee}{\underline C}_{A|R}
 \longrightarrow\widetilde{\underline C}_{A|R}
\]
为伴随层函子.

现假定 $F$ 加性.
断言对每个 $q>0$, $\beta^{\#}\underline H^q_{A|S}(F)$ 都是常层.
事实上, 若 $P$ 为 $R$ 上的多项式代数,
则 $P=R\otimes_SP'$, 其中 $P'$ 为 $S$ 上的多项式代数, 且
\[
 [\beta^{\#}\underline H^q_{A|S}(F)](P)
 =H^q(P|S,F)
 =H^q(R\otimes_SP'|S,F)
 \longrightarrow H^q(R|S,F)
\]
由 (b) 在 $q>0$ 时为同构, 在 $q=0$ 时为满射.
因此由 \SGARefLink{sga7.sga7i.evi.n3-3}{3.3}(c), 到相应常层的典范映射
\[
 \beta^{\#}\underline H^q_{A|S}(F)
 \longrightarrow H^q(R|S,F)
\]
在 $q>0$ 时为同构, 在 $q=0$ 时为满射.
故只要 $p,q>0$ 就有 $E^{p,q}=0$, 谱序列给出正合序列
\[
 (*)\quad
 0\longrightarrow E^{1,0}\longrightarrow H^1
 \longrightarrow E^{0,1}\longrightarrow E^{2,0}
 \longrightarrow H^2\longrightarrow E^{0,2}
 \longrightarrow E^{3,0}\longrightarrow H^3\longrightarrow\cdots.
\]
此处
\[
 E^{0,q}=[\beta^{\#}\underline H^q_{A|S}(F)](A)
 \xrightarrow{\ \sim\ }H^q(R|S,F)\quad(q>0),
 \qquad
 E^{p,0}=H^p(A|R,\widetilde\beta F).
\]
另一方面, 正合序列
\[
 0\longrightarrow F_R\longrightarrow\widetilde\beta F
 \longrightarrow F(R)\longrightarrow0
\]
给出正合图
\begin{center}
\scriptsize
\begin{tikzcd}[column sep=small,row sep=large]
0 \arrow[r]
& H^0(A|R,F_R) \arrow[r]
& H^0(A|R,\widetilde\beta F) \arrow[r] \arrow[d,equal]
& H^0(A|R,F(R)) \arrow[r] \arrow[d,equal]
& H^1(A|R,F_R) \arrow[r]
& H^1(A|R,\widetilde\beta F) \arrow[r]
& 0\\
&& H^0(A|S,F) & H^0(R|S,F) &&&
\end{tikzcd}
\end{center}
原文将此图记为 $(**)$. 结合 $(*)$ 与 $(**)$ 即得所需的正合序列. \hfill$\square$

% END INCLUDED SGA7I_English_current_buildable_source_20260823/components/30_expose_VI_proposition_3_7.tex

% BEGIN INCLUDED SGA7I_English_current_buildable_source_20260823/components/31_expose_VI_corollaries_3_8_and_3_9.tex
%%%%%%%%%%  scan idx 92  |  folio 81  %%%%%%%%%%

\underline{推论 \SGARefTarget{sga7.sga7i.evi.n3-8}{3.8}.}

\begin{enumerate}[label=(\alph*)]
\item 设 $A$ 和 $B$ 为 $R$-代数, 满足
\[
 \operatorname{Tor}_i^R(A,B)=0\qquad(0<i\leq n).
\]
则对每个加性层
$F\in\operatorname{Ob}(\widetilde{\underline C}_{A\otimes_RB|R})$, 典范映射
\[
 H^i(A\otimes_RB|R,F)
 \longrightarrow H^i(A|R,F)\oplus H^i(B|R,F)
\]
对每个 $i\leq n$ 都是同构.

\item 设 $A$ 为 $R$-代数, 满足对每个 $i>0$ 有
$\operatorname{Tor}_i^R(A,A)=0$.
若映射 $A\otimes_RA\longrightarrow A$ 为同构,
则对每个加性层 $F\in\operatorname{Ob}(\widetilde{\underline C}_{A|R})$, 有
\[
 H^i(A|R,F)=0
\]
对每个 $i$ 成立.
\end{enumerate}

\underline{证明.}

(a) 设 $F\in\operatorname{Ob}(\widetilde{\underline C}_{A\otimes_RB|R})$ 加性.
对 $R\longrightarrow A\longrightarrow A\otimes_RB$ 应用命题
\SGARefLink{sga7.sga7i.evi.n3-7}{3.7}(c), 得到正合序列
\[
 (*)\quad\cdots\longrightarrow
 H^i(A\otimes_RB|A,F_A)
 \longrightarrow H^i(A\otimes_RB|R,F)
 \longrightarrow H^i(A|R,F)\longrightarrow\cdots.
\]
考虑交换图
\[
\begin{tikzcd}[column sep=huge,row sep=huge]
H^i(A\otimes_RB|A,F_A) \arrow[r] \arrow[d]
& H^i(A\otimes_RB|R,F) \arrow[d]\\
H^i(B|R,\widetilde A F_A) \arrow[r,"\approx"]
& H^i(B|R,F).
\end{tikzcd}
\]
此处
\[
 A:\underline C_{B|R}\longrightarrow
 \underline C_{A\otimes_RB|A},
 \qquad X\longmapsto A\otimes_RX,
\]
下方映射由 $\widetilde A F_A\longrightarrow F$ 诱导;
因 $F$ 加性, 此映射为同构.
由函子性, 图的交换性归结为 $i=0$ 的情形.

若对 $0<i\leq n$ 有 $\operatorname{Tor}_i^R(A,B)=0$,
则命题 \SGARefLink{sga7.sga7i.evi.n3-7}{3.7}(a)
表明左侧竖直映射对 $i\leq n$ 为同构. 从而复合
\[
 H^i(A\otimes_RB|A,F_A)
 \longrightarrow H^i(A\otimes_RB|R,F)
 \longrightarrow H^i(B|R,F)
\]
及由对称性得到的复合

%%%%%%%%%%  scan idx 93  |  folio 82  %%%%%%%%%%

\[
 H^i(A\otimes_RB|B,F_B)
 \longrightarrow H^i(A\otimes_RB|R,F)
 \longrightarrow H^i(A|R,F)
\]
在 $i\leq n$ 时均为同构.
因此 $(*)$ 在 $i\leq n$ 的范围内分裂为分裂短正合序列
\[
 0\longrightarrow H^i(A\otimes_RB|A,F_A)
 \longrightarrow H^i(A\otimes_RB|R,F)
 \longrightarrow H^i(A|R,F)\longrightarrow0
\]
故
\[
 H^i(A\otimes_RB|R,F)
 \longrightarrow H^i(A|R,F)\oplus H^i(B|R,F)
\]
在此范围内为同构.

(b) 由 (a), 映射
\[
 H^i(A\otimes_RA|R,F)
 \longrightarrow H^i(A|R,F)\oplus H^i(A|R,F)
\]
对每个 $i$ 均为同构.
另一方面, 假设中的同构 $A\otimes_RA\longrightarrow A$ 诱导同构
\[
 H^i(A|R,F)\longrightarrow H^i(A\otimes_RA|R,F).
\]
其复合为对角映射
\[
 H^i(A|R,F)\longrightarrow
 H^i(A|R,F)\oplus H^i(A|R,F).
\]
因它为同构, 对每个 $i$ 有 $H^i(A|R,F)=0$. \hfill$\square$

\underline{推论 \SGARefTarget{sga7.sga7i.evi.n3-9}{3.9}.}
设 $S$ 为 $R$-代数 $A$ 的乘法子集. 假定
\[
 F\in\operatorname{Ob}
 (\widetilde{\underline C}_{S^{-1}A|R})
\]
加性, 且 $F_A$ 也加性. 则典范映射
\[
 H^i(S^{-1}A|R,F)\longrightarrow H^i(A|R,F)
\]
对每个 $i$ 均为同构.
此处 $F_A\in\operatorname{Ob}(\widetilde{\underline C}_{S^{-1}A|A})$
对每个 $X\in\operatorname{Ob}(\underline C_{S^{-1}A|A})$ 定义为
\[
 F_A(X)=\ker\bigl(F(X)\longrightarrow F(A)\bigr)
\]

\underline{证明.} 对 $R\longrightarrow A\longrightarrow S^{-1}A$
应用命题 \SGARefLink{sga7.sga7i.evi.n3-7}{3.7}(c), 得到正合序列
\[
 \cdots\longrightarrow H^n(S^{-1}A|A,F_A)
 \longrightarrow H^n(S^{-1}A|R,F)
 \longrightarrow H^n(A|R,F)\longrightarrow\cdots.
\]
映射
\[
 S^{-1}A\otimes_AS^{-1}A\longrightarrow S^{-1}A
\]
为同构, 且由假设 $F_A$ 加性.
因此由推论 \SGARefLink{sga7.sga7i.evi.n3-8}{3.8}(b),
对每个 $n$ 有 $H^n(S^{-1}A|A,F_A)=0$.
因而正合序列的中间映射对每个 $n$ 均为同构. \hfill$\square$

% END INCLUDED SGA7I_English_current_buildable_source_20260823/components/31_expose_VI_corollaries_3_8_and_3_9.tex

% BEGIN INCLUDED SGA7I_English_current_buildable_source_20260823/components/32_expose_VI_definition_and_lemma_3_10.tex
%%%%%%%%%%  scan idx 94  |  folio 83  %%%%%%%%%%

设 $A$ 为 $R$-代数. $A$-模 $M$ 定义加性层
\[
 \operatorname{Der}_R(-,M)
 \in\operatorname{Ob}(\widetilde{\underline C}_{A|R}),
 \qquad X\longmapsto\operatorname{Der}_R(X,M).
\]
作如下定义.

\underline{定义 \SGARefTarget{sga7.sga7i.evi.n3-10}{3.10}.}
设 $A$ 为 $R$-代数. 对每个 $A$-模 $M$, 置
\[
 D^i(A|R,M)=H^i(A|R,\operatorname{Der}_R(-,M)).
\]

现计算低次群, 即 $i=1,2$ 时的 $D^i(A|R,M)$.
在 $\underline C_{A|R}$ 中取满射箭头 $P\longrightarrow A$,
其中 $P$ 为 $R$ 上的多项式代数, 并写成
\[
 0\longrightarrow I\longrightarrow P\longrightarrow A
 \longrightarrow0.
\]

\underline{引理 \SGARefTarget{sga7.sga7i.evi.n3-10.occ-2}{3.10}.}
\begin{align*}
 D^1(A|R,M)
 &=\operatorname{Coker}\bigl(
   \operatorname{Der}_R(P,M)\longrightarrow
   \operatorname{Hom}_P(I,M)\bigr),\\
 D^2(A|R,M)
 &=H^1\bigl(A|R,
   \underline H^1_{A|R}(\operatorname{Der}_R(-,M))\bigr).
\end{align*}

\underline{证明.} 考虑单纯代数
\[
 (*)\qquad
 \cdots\ P_A^{(2)}
 \mathrel{\substack{\longrightarrow\\[-.6ex]\longrightarrow\\[-.6ex]\longrightarrow}}
 P_A^{(1)}
 \mathrel{\substack{\longrightarrow\\[-.6ex]\longrightarrow}}
 P_A^{(0)},
\]
其中
\[
 P_A^{(n)}=
 \overbrace{P\times_A\cdots\times_AP}^{n+1}.
\]
对每个 $n$, 有正合序列
\[
 0\longrightarrow I^{(n)}\longrightarrow P_A^{(n)}
 \longrightarrow A\longrightarrow0,
 \qquad
 I^{(n)}=\overbrace{I\times\cdots\times I}^{n+1}.
\]
置 $\overline I=I/I^2$, 并以 $\Delta$ 表示对角映射.
易于验证, 对每个 $n$, 序列
\[
 0\longrightarrow
 \operatorname{Hom}_A
 \bigl(\overline I^{(n)}/\Delta(\overline I),M\bigr)
 \longrightarrow\operatorname{Der}_R(P_A^{(n)},M)
 \xrightarrow{\operatorname{Der}(\Delta,M)}
 \operatorname{Der}_R(P,M)\longrightarrow0
 \tag{a}
\]
正合. 还有正合序列
\[
 0\longrightarrow
 \operatorname{Hom}_A
 \bigl(\overline I^{(n)}/\Delta(\overline I),M\bigr)
 \longrightarrow\operatorname{Hom}_A(\overline I^{(n)},M)
 \xrightarrow{\operatorname{Hom}(\Delta,M)}
 \operatorname{Hom}_A(\overline I,M)\longrightarrow0.
 \tag{b}
\]
序列 (a),(b) 给出相应的单纯 $A$-模的正合序列.
因 $\{I\longrightarrow0\}$ 及 $\{I\longrightarrow I\}$ 有截面,
\[
 H^i(\{I\longrightarrow0\},\operatorname{Hom}_A(-,M))
 =H^i(\{I\longrightarrow I\},\operatorname{Hom}_A(-,M))=0
 \qquad(i>0).
\]

%%%%%%%%%%  scan idx 95  |  folio 84  %%%%%%%%%%

由此得到
\[
\left\{
\begin{aligned}
0\longrightarrow\operatorname{Der}_R(A,M)
&\longrightarrow\operatorname{Der}_R(P,M)
\longrightarrow\operatorname{Hom}_A(\overline I,M)\\
&\longrightarrow
H^1(\{P\longrightarrow A\},\operatorname{Der}_R(-,M))
\longrightarrow0
&&\text{正合},\\
H^i(\{P\longrightarrow A\},\operatorname{Der}_R(-,M))
&=0 &&(i>1).
\end{aligned}
\right.
\]
因此由 \SGARefLink{sga7.sga7i.evi.n3-1.item-b}{3.1}(b),
\begin{align*}
D^1(A|R,M)
 &=\check H^1(A|R,\operatorname{Der}_R(-,M))\\
 &=H^1(\{P\longrightarrow A\},\operatorname{Der}_R(-,M))\\
 &=\operatorname{Coker}\bigl(
   \operatorname{Der}_R(P,M)\longrightarrow
   \operatorname{Hom}_A(\overline I,M)\bigr)\\
 &=\operatorname{Coker}\bigl(
   \operatorname{Der}_R(P,M)\longrightarrow
   \operatorname{Hom}_P(I,M)\bigr),\\
D^2(A|R,M)
 &=H^1\bigl(A|R,
   \underline H^1_{A|R}(\operatorname{Der}_R(-,M))\bigr).
\end{align*}
\hfill$\square$

取自由 $P$-模 $F$ 及 $P$-线性映射 $\gamma:F\longrightarrow P$,
使 $\gamma(F)=I$. 与 $\gamma$ 相伴的是一个 Koszul 复形,
稍微滥用记号, 将其记为 $K_I$.

% END INCLUDED SGA7I_English_current_buildable_source_20260823/components/32_expose_VI_definition_and_lemma_3_10.tex
